\documentclass[11pt,letterpaper]{article}
\usepackage[T1]{fontenc}
\usepackage[utf8]{inputenc}
\usepackage{lmodern}
\usepackage[margin=1in,headheight=15pt]{geometry}
\usepackage{amsmath,amssymb,amsthm,mathtools,mathrsfs}
\usepackage{microtype,booktabs,array,enumitem}
\usepackage[dvipsnames]{xcolor}
\usepackage{fancyhdr,aliascnt}
\usepackage{xurl}
\usepackage[colorlinks=true,linkcolor=MidnightBlue,citecolor=MidnightBlue,urlcolor=MidnightBlue,breaklinks=true]{hyperref}
\usepackage[nameinlink,noabbrev]{cleveref}
\hypersetup{%
  pdftitle={Nonabelian Support Obstructions: A Matrix Cousin Criterion and the Surcomplex Riemann-Hilbert Problem},
  pdfauthor={Merged research manuscript},
  pdfsubject={Positive Hahn matrix gluing, polar supports, and inverse monodromy}}
\pagestyle{fancy}
\fancyhf{}
\fancyhead[L]{\small Nonabelian support obstructions}
\fancyhead[R]{\small Matrix Cousin and Riemann--Hilbert}
\fancyfoot[C]{\thepage}
\setlength{\parskip}{0.28em}
\setlength{\parindent}{1.2em}
\setlist{nosep,leftmargin=2em}
\emergencystretch=2em

\newtheorem{theorem}{Theorem}[section]
\newaliascnt{lemma}{theorem}
\newtheorem{lemma}[lemma]{Lemma}
\aliascntresetthe{lemma}
\newaliascnt{proposition}{theorem}
\newtheorem{proposition}[proposition]{Proposition}
\aliascntresetthe{proposition}
\newaliascnt{corollary}{theorem}
\newtheorem{corollary}[corollary]{Corollary}
\aliascntresetthe{corollary}
\theoremstyle{definition}
\newaliascnt{definition}{theorem}
\newtheorem{definition}[definition]{Definition}
\aliascntresetthe{definition}
\newaliascnt{example}{theorem}
\newtheorem{example}[example]{Example}
\aliascntresetthe{example}
\theoremstyle{remark}
\newaliascnt{remark}{theorem}
\newtheorem{remark}[remark]{Remark}
\aliascntresetthe{remark}

\newcommand{\C}{\mathbb C}
\newcommand{\R}{\mathbb R}
\newcommand{\Q}{\mathbb Q}
\newcommand{\Z}{\mathbb Z}
\newcommand{\N}{\mathbb N}
\newcommand{\No}{\mathbf{No}}
\newcommand{\SC}{\mathbf{SC}}
\newcommand{\OO}{\mathcal O}
\newcommand{\HH}{\mathscr H_\Gamma}
\newcommand{\II}{\mathscr I_\Gamma}
\newcommand{\JJ}{\mathscr J_\Gamma}
\newcommand{\GG}{\mathscr G_{r,\Gamma}}
\newcommand{\K}{K_\Gamma}
\newcommand{\mm}{\mathfrak m_\Gamma}
\newcommand{\LL}{\mathcal L}
\newcommand{\Mat}{\operatorname{Mat}}
\newcommand{\GL}{\operatorname{GL}}
\newcommand{\SL}{\operatorname{SL}}
\newcommand{\supp}{\operatorname{supp}}
\newcommand{\red}{\operatorname{red}}
\newcommand{\Pol}{\operatorname{Pol}}
\newcommand{\Reg}{\operatorname{Reg}}
\newcommand{\Hol}{\operatorname{Hol}}
\newcommand{\Log}{\operatorname{Log}}
\newcommand{\Res}{\operatorname{Res}}
\newcommand{\tr}{\operatorname{tr}}
\newcommand{\val}{\operatorname{v}}
\newcommand{\im}{\operatorname{im}}
\newcommand{\id}{\operatorname{id}}
\newcommand{\Endsh}{\mathcal End}
\newcommand{\PP}{\mathcal P}
\newcommand{\EU}{\mathcal E}
\newcommand{\coeff}[1]{[t^{#1}]}
\newcommand{\angles}[1]{\langle #1\rangle_{\N}}
\newcolumntype{L}[1]{>{\raggedright\arraybackslash}p{#1}}
\numberwithin{equation}{section}

\begin{document}
\hypersetup{pageanchor=false}
\begin{titlepage}
\centering
\vspace*{0.4cm}
{\large\scshape A merged research continuation of the Surreal repository\par}
\vspace{0.7cm}
{\Huge\bfseries Nonabelian Support\par Obstructions\par}
\vspace{0.5cm}
{\LARGE A Matrix Cousin Criterion\par and the Surcomplex\par Riemann--Hilbert Problem\par}
\vspace{0.7cm}
{\large September 21, 2026\par}
\vspace{0.55cm}
\begin{minipage}{0.94\textwidth}
\begin{abstract}
We study one coefficient category from two sides. Its objects are matrices
congruent to the identity modulo positive Hahn exponents,
$I_r+\Mat_r(\JJ)$, over a set-sized and possibly non-Archimedean ordered
value group $\Gamma$, with ordinary holomorphic coefficient functions on a
single fixed ordinary complex domain; when $\Gamma\subseteq\No$, the
substitution $t^\gamma=\omega^{-\gamma}$ gives the surcomplex reading.

Part~I solves the multiplicative gluing problem. On a puncture-and-disk cover
of a connected open Riemann surface we prove an exact criterion: factor each
transition matrix \emph{independently} into a normalized polar factor and a
factor holomorphic on the full disk, and test whether the union of the
supports of the \emph{polar} factors is well ordered. The criterion is
invariant under arbitrary local holomorphic right gauges and under ordered
value-group extension, and it produces genuinely nonabelian counterexamples:
a rank-three family whose determinant is one, whose only raw singular
exponent is $1$, and whose abelianized problem splits, yet whose integral
vector bundle is nontrivial; a repair of that family by a term which
\emph{worsens} its raw singular support; and, for every $r\ge3$, an example
trivial modulo the last central subgroup of the upper-unitriangular group but
nonsplit in full positive rank $r$. We also prove support-controlled Stein
splitting and deformation rigidity, with a cyclic-versus-noncyclic dichotomy.

Part~II solves the inverse monodromy problem in the same category. A
near-identity representation of $\pi_1(\C\setminus D)$ is the monodromy of a
positive globally common-domain Hahn connection exactly when the union over
punctures of the supports of $\rho(\ell_d)-I_r$ is well ordered; the realizer
may then be taken coefficientwise logarithmic, and a fixed linear
Mittag--Leffler period section selects a unique normalized realizer. We prove
framed gauge classification, entire extension of logarithmic comparison
gauges, a positive Frobenius form, a determinant-one refinement, calibrating
examples on both sides of the criterion, an explicit mixed commutator
correction in the value group $\Q+\Q\omega$, and the impossibility of a
continuous linear period right inverse for infinitely many punctures.

\medskip
\noindent\textbf{The two headline criteria are different criteria on
different objects.} Part~I tests the supports of normalized polar factors, not
the raw transition supports; Part~II tests the raw meridian monodromy
supports. \Cref{nab:sec:twocriteria} states which object each is a criterion
on, and \cref{nab:sec:rankthree} contains the example that makes the
distinction unavoidable.
\end{abstract}
\end{minipage}
\vfill
\begin{minipage}{0.94\textwidth}\small
\textbf{Status and scope.} This report merges two independently written
manuscripts; see \cref{nab:app:repro}. The exact global criteria and the
explicit obstruction families are proposed new results in the specified
Hahn-coherent categories. The local factorization mechanisms, the Neumann
support calculus and every ordinary analytic input are established
mathematics and are credited as such. The literature checks were targeted and
certify no priority; no named published conjecture is claimed solved. Nothing
here is refereed or machine checked. Nontriviality in Part~I is proved over
the \emph{integral} coefficient sheaf only, and Part~II assumes positive
support and near-identity monodromy throughout: neither theorem is restated
over a coefficient field in which negative-valuation changes of frame are
admitted, and neither is the unrestricted classical Cousin or
Riemann--Hilbert theorem. All infinite algebra is by well-ordered supports and
finite coefficient multiplicity, never by convergence in the fine surreal
topology.
\end{minipage}
\end{titlepage}
\hypersetup{pageanchor=true}
\pagenumbering{roman}
\tableofcontents
\newpage
\pagenumbering{arabic}

\section{Two questions, one category}\label{nab:sec:intro}

\subsection{The shared setting}

Let $\Gamma$ be a set-sized totally ordered abelian group and let
$t^\Gamma$ be its monomial group. A coefficient of our objects is an ordinary
holomorphic function on a \emph{single fixed} ordinary complex domain, and an
object is a Hahn series in $t$ with well-ordered support whose coefficients
all live on that same domain. The matrices we glue and the matrices we realize
as monodromy both lie in
\[
 I_r+\Mat_r\bigl(\text{positive-support part}\bigr),
\]
so they are congruent to the identity modulo positive exponents. When
$\Gamma$ is an ordered additive subgroup of the surreal numbers $\No$, the
Conway normal form substitution $t^\gamma=\omega^{-\gamma}$ interprets
everything inside $\SC=\No[i]$. That interpretation is a reading of a Hahn
model; it is never an assertion that a fine-topological surcomplex manifold or
a fine-continuous surcomplex path is involved.

Two problems are asked in this one category.

\paragraph{Part I: multiplicative gluing (a matrix Cousin problem).}
Let $X$ be a connected open Riemann surface, let $A\subset X$ be closed and
discrete, let $(D_a,u_a)$ be pairwise disjoint coordinate disks, and put
$U_0=X\setminus A$, $W_a=D_a\setminus\{a\}$. Given, \emph{separately for each
$a$}, a positive matrix $G_a$ on $W_a$, when are there positive matrices
$B_0$ on $U_0$ and $B_a$ on $D_a$ with
\begin{equation}\label{nab:eq:splitintro}
 G_a=B_0B_a^{-1}\quad\text{on }W_a?
\end{equation}
Each local equation is harmless. The difficulty is that the entries of the
\emph{one global matrix} $B_0$ must have a legitimate, that is well ordered,
Hahn support.

\paragraph{Part II: inverse monodromy (a Riemann--Hilbert problem).}
Let $D\subset\C$ be closed and discrete, $U=\C\setminus D$, and let
$\rho:\pi_1(U,p)\to I_r+\Mat_r(\mm)$ be a near-identity representation. When
is $\rho$ the monodromy of a connection $dY=\Omega Y$ whose coefficients are
ordinary holomorphic matrix one-forms on all of $U$ and whose Hahn support is
a single well-ordered positive set?

These are different questions sharing a coefficient category and an ordinary
Mittag--Leffler input applied one exponent at a time. Part~I first normalizes
each transition matrix; Part~II corrects the lower-weight monodromy prediction
by ordinary periods. Mixed products matter in both constructions, but the
support obstruction in Part~II can already occur for an abelian
representation. In both problems a certain \emph{union of Hahn supports must
be well ordered}. The objects in those unions must be distinguished.

\subsection{The two criteria, and why they are two}

\begin{center}
\fbox{\begin{minipage}{0.93\textwidth}
\textbf{Criterion P} (Part~I, \cref{nab:thm:criterion}). Gluing
\eqref{nab:eq:splitintro} is solvable if and only if
\[
 \EU(G)=\bigcup_{a\in A}\supp\bigl(\Pol(G_a)-I_r\bigr)
 \quad\text{is well ordered,}
\]
where $\Pol(G_a)$ is the normalized polar factor produced by an
\emph{independent local factorization} of each transition matrix
$G_a=\Pol(G_a)\Reg(G_a)$. This is a criterion on the polar factors, not on
the transition matrices. Well-ordering of the full raw union
$\bigcup_a\supp(G_a-I_r)$ is \emph{sufficient but not necessary}.
Well-ordering of the raw \emph{singular-coefficient} support is neither
necessary nor sufficient (\cref{nab:sec:rankthree}).

\medskip
\textbf{Criterion M} (Part~II, \cref{nab:thm:RH}). A near-identity
representation $\rho$ is realizable if and only if
\[
 E_\rho=\bigcup_{d\in D}\supp\bigl(\rho(\ell_d)-I_r\bigr)
 \quad\text{is well ordered,}
\]
the union being over the matrices themselves, for a compatible system of
based meridians $\ell_d$ freely generating $\pi_1(U,p)$. No factorization is
applied first, and none may be: this is a criterion on the monodromy
matrices, and it is basis- and basepoint-independent
(\cref{nab:prop:intrinsic}).

\medskip
A reader who carries one ``well-ordered support criterion'' away from this
report will test the wrong object in one of the two problems.
\end{minipage}}
\end{center}

\medskip
The distinction is not pedantic. \Cref{nab:sec:rankthree} contains a
rank-three family over $\Gamma=\Q$ whose raw singular support is the single
exponent $\{1\}$ and which nevertheless cannot be glued, because the
normalized polar factors are forced to carry the descending set
$\{1+1/n\}$; and it contains the opposite example, in which \emph{adding} a
term that introduces exactly that descending set among the raw poles
\emph{removes} the obstruction. So in Part~I the raw singular support is
neither sufficient nor necessary, and the factorization must be performed
first. In Part~II, by contrast, well-ordering of the raw monodromy support is
the invariant condition, and applying a factorization to $\rho(\ell_d)$ before testing
would be a mistake of the opposite kind.

\subsection{What each part contributes}

Part~I owns the machinery the two parts share: the integral common-domain
Hahn sheaf on an ordinary base, the noncommutative support calculus, the
normalized local polar factorization with right-gauge invariance, invariance
under ordered value-group extension, the recovery of the known rank-one
scalar test, and the Stein branch. Its own new content is the exact criterion
$\EU(G)$, the rank-three cross-term obstruction with its repair and its
logarithm diagnostic, the arbitrary-depth central hierarchy, the resulting
determinant-one integral vector bundles that are trivial on every bounded
plane domain, the explicit localization boundary, and the
cyclic-versus-noncyclic rigidity dichotomy.

Part~II reads as the connection-and-monodromy application of that machinery,
with its own new content: the three-way equivalence of
\cref{nab:thm:RH} including \emph{necessity}, which comes from a uniform
support bound on holonomy valid for every group element; the fixed
complex-linear period section that selects a unique normalized realizer; the
framed gauge classification and the entire extension of logarithmic
comparison gauges; the positive Frobenius form in which no spectral resonance
operator is ever inverted; the $\SL_r$ refinement; three calibrating examples
that pin the criterion from both sides; the explicit first nonabelian
correction with its $\Q+\Q\omega$ specialization; and the proof that an
unrestricted infinite-puncture linear period right inverse cannot also be
continuous.

\subsection{Relation to reports already in the collection}

This report is a continuation of the repository
\texttt{VladimirReshetnikov/Surreal}, reviewed at commit
\begin{center}\small\ttfamily
39f2be6667ade51bca2b45daa47e289d69c09764,
\end{center}
inspected September 21, 2026. Two companion reports are leaned on and are
\emph{cited rather than reproved}.

\emph{Global Support Obstructions: Moving Divisors, Mittag--Leffler, and a
Picard Dichotomy}, at \texttt{docs/surcomplex/global-divisors/article.tex}
\cite{RepoGlobal}, proves the scalar support criteria and the Picard-group
dichotomy for this workspace, and its subsection ``Further precise
directions'' asks for exactly the present work: ``the finite-rank
vector-bundle theory: matrix-valued positive transition functions no longer
reduce to an abelian logarithm identity, so a noncommutative support
recursion must replace the rank-one argument'', and the same package over
general ordinary open Riemann surfaces. Part~I answers both. Its rank-one
specialization \emph{recovers} that report's exact Cousin criterion, its
Theorem~13.1, stated there in terms of the singular support $\Sigma(b)$; we
cite that theorem as the scalar case rather than reprove it
(\cref{nab:prop:rankone}). Likewise the Picard-group vanishing dichotomy of
that report --- vanishing exactly when $\Gamma=\{0\}$ or $\Gamma=\Z\delta$ ---
is the rank-one shadow of our Stein dichotomy
(\cref{nab:thm:dichotomy}), driven by the same universal well-ordered set
$\{\delta,2\delta,\ldots\}$.

\emph{Differential Algebra and Differential Equations over the Surreal and
Surcomplex Numbers}, at
\texttt{docs/surcomplex/differential-equations/article.tex}
\cite{RepoDiff}, contains the \emph{forward} map that Part~II inverts: its
theorem ``Monodromy valued in $\GL_d(K_\Gamma)$'' constructs $M_\ell$ in
$\GL_d(\K)$ from a coherent connection, records that the coefficient of
$M_\ell$ at exponent $0$ is the ordinary monodromy matrix and that all other
exponents lie in $S^*\setminus\{0\}$, and then states the categorical
boundary we adopt verbatim rather than re-argue: ``No fine-continuous
surcomplex path is asserted: the loop is an ordinary complex path,
continuation is coefficientwise, and the result is a matrix of Hahn
constants.'' \Cref{nab:thm:RH} is the converse of that map for near-identity
representations, together with the exact condition on which ones arise.

\subsection{What is and is not claimed as new}

The Hahn support calculus is classical \cite{Neumann,Poonen}. The analytic
interpretation of restricted ordinary functions over surreal workspaces is
classical \cite{vdDE}, and the broader landscape is surveyed in
\cite{MantovaMatusinski}. Algebraic Birkhoff factorization in complete
filtered noncommutative algebras, including its recursive form and the
necessity of commutator corrections, is established \cite{EFGK,EFMP}, as are
the ordinary Laurent, Mittag--Leffler, Behnke--Stein and Cartan inputs
\cite{BS,Cartan,Schmitt}, iterated path integrals \cite{Chen}, and the
classical regular-singular and Riemann--Hilbert background
\cite{Deligne,Ilyashenko}. We reprove only what we need with explicit
arbitrary-ordered-group support control, and we claim no new general
factorization principle and no new ordinary analytic theorem.

We do not classify bundles or connections beyond the stated questions. Part~I
decides triviality on a puncture-and-disk cover for arbitrary finite rank and
proves a sufficient theorem for arbitrary Stein covers; Part~II inverts
holonomy for near-identity representations on a punctured plane. The
boundaries in \cref{nab:sec:scope} are part of the theorem statements, not
qualifications added afterwards.

\section{The coefficient category}\label{nab:sec:category}

\subsection{One set-sized workspace}

For a complex algebra $R$, the Hahn algebra $R((t^\Gamma))$ consists of formal
sums $F=\sum_{\gamma\in S}f_\gamma t^\gamma$ with $f_\gamma\in R$ whose
nonzero support $S$ is well ordered in the given order; addition is
coefficientwise and multiplication is convolution. For matrices the support
is the union of the supports of the finitely many entries, and an exponent
occurs exactly when the corresponding coefficient matrix is not identically
zero. The Hahn field of constants is
\[
 \K=\C((t^\Gamma)),\qquad
 \mm=\{a\in\K:\val(a)>0\},
\]
where $\val$ of a nonzero series is its least exponent and $\val(0)=+\infty$.
Divisibility of
$\Gamma$ and algebraic closedness of $\K$ are never used: we extract no roots
and take no algebraic closure.

When $\Gamma\subseteq\No$ is an ordered additive subgroup, Conway normal forms
give the interpretation
\begin{equation}\label{nab:eq:surrealembedding}
 \C((t^\Gamma))\longrightarrow\No[i]=\SC,
 \qquad \sum c_\gamma t^\gamma\longmapsto\sum c_\gamma\omega^{-\gamma},
\end{equation}
with real and imaginary parts interpreted separately
\cite{vdDE,Poonen,Bournez}. All supports here are sets; no sum indexed by a
proper class, no proper-class fundamental group and no proper-class
coefficient ring occurs anywhere.

All ordinary manifolds and Riemann surfaces carry their usual Hausdorff,
second-countable topology. For a connected ordinary complex domain $V$ set
\begin{align}
 \HH(V)&=\OO(V)((t^\Gamma)),\label{nab:eq:H}\\
 \II(V)&=\{F\in\HH(V):\supp F\subseteq\Gamma_{\geq0}\},\label{nab:eq:I}\\
 \JJ(V)&=\{F\in\HH(V):\supp F\subseteq\Gamma_{>0}\}.\label{nab:eq:J}
\end{align}
Every coefficient must be holomorphic on the \emph{same fixed domain} $V$; we
never substitute a family of convergent germs with unrelated radii, and
$\HH(V)$ is in general \emph{not} the algebraic tensor product
$\K\otimes_\C\OO(V)$, since a legitimate family may involve infinitely many
linearly independent coefficient functions. For disconnected open sets,
sections are specified component by component. Writing
$\Omega^1_{\mathrm{hol}}(V)$ for ordinary holomorphic one-forms, the same
notation with one-form coefficients denotes one-form families. There is a
reduction map
\begin{equation}\label{nab:eq:reduction}
 \red:\II\longrightarrow\OO,\qquad
 \sum_{\gamma\geq0}f_\gamma t^\gamma\longmapsto f_0,
 \qquad\ker\red=\JJ,
\end{equation}
and the positive matrix group is
\begin{equation}\label{nab:eq:G}
 \GG(V)=I_r+\Mat_r(\JJ(V)).
\end{equation}
Its elements are automatically two-sided invertible by
\cref{nab:lem:inverse}. Throughout, ``positive'' refers to \emph{Hahn
exponents only}: never to positivity of a Hermitian matrix, never to
positivity of entries, and never to positive-definiteness. Elements of
$I_r+\Mat_r(\mm)$ need not be algebraically unipotent.

\subsection{Two standing conventions that differ between the two parts}
\label{nab:sub:conventions}

The two parts of this report were developed independently and do not make
identical structural assumptions. Both assumptions are correct as stated and
neither is weakened here.

Part~I \emph{proves and uses} the sheaf property of these assignments on the
ordinary base. The point is a support point. A nonzero ordinary holomorphic
coefficient on a connected domain cannot vanish on a nonempty open subset, so
restriction between nonempty connected ordinary domains preserves support. If
compatible local sections are given on a cover of a connected domain, the
identity theorem along chains of intersecting patches shows that their
coefficient functions glue and that the global support equals the support seen
on any one nonempty connected patch; that support is well ordered. For
disconnected opens one takes the product over connected components. In
particular, if a section on $X\setminus A$ extends separately across each
point of a discrete $A$, it extends \emph{with its existing global support}:
no new exponent can appear only at an isolated ordinary point. This
observation is used whenever a matrix quotient is extended across a puncture.

Part~II deliberately does \emph{not} assume that the global-support assignment
is a sheaf for arbitrary covers, and none of its proofs uses that. All its
gluing is done directly, by descending families already supported in one
specified monoid, and its only analytic input is the ordinary
Mittag--Leffler theorem applied one exponent at a time. Consequently no
application of a Cartan-type theorem to an infinite-dimensional Hahn
coefficient sheaf is hidden in Part~II, and no companion claim about Hahn
Noetherianity or a Hahn Nullstellensatz is relied upon anywhere in this
report.

\subsection{Comparison of names across the two parts}
\label{nab:sub:names}

The following table fixes one vocabulary. Where the two source manuscripts
used different names for the same object, the merged name is given first;
where they used the same word for genuinely different objects, the objects are
separated.

\begin{center}\small
\begin{tabular}{@{}L{0.26\textwidth}L{0.33\textwidth}L{0.33\textwidth}@{}}
\toprule
Merged name & Part~I usage & Part~II usage\\
\midrule
$\HH(V),\II(V),\JJ(V)$ & full, integral and positive common-domain sheaves;
sheaf property proved on the ordinary base &
$\JJ(V)$ is the positive part of the common-domain family; no sheaf property
assumed\\[0.35em]
$\GG(V)=I_r+\Mat_r(\JJ(V))$ & the positive matrix group being split &
the group of positive gauges; framed gauges have value $I_r$ at the base
point, and $\rho$ takes values in the constant subgroup\\[0.35em]
$\mm$ & not used; positivity is imposed sheaf-wise &
positive-valuation elements of $\K$; $\rho$ takes values in
$I_r+\Mat_r(\mm)$\\[0.35em]
Well-ordered union of supports & taken over normalized \emph{polar factors}
$\Pol(G_a)-I_r$: well-ordering of $\EU(G)$ is invariant &
taken over the \emph{monodromy matrices} $\rho(\ell_d)-I_r$: well-ordering of
$E_\rho$ is invariant\\[0.35em]
Local positive factorization & $G=\Pol(G)\Reg(G)$ on a punctured disk,
normalized Birkhoff type &
$Y=H\exp(R\Log z)$ on a punctured disk, positive Frobenius type\\[0.35em]
$\Sigma(\cdot)$ & the rank-one logarithmic singular support, used only for
the scalar comparison & not used\\[0.35em]
Strong summability & implicit in Hahn summability of finite expressions &
a named condition on set-indexed families; distinguished sharply from
valuation-topology convergence\\
\bottomrule
\end{tabular}
\end{center}

\begin{remark}[Resonance does not occur in this report]
\label{nab:rem:resonance}
Neither part imposes, or needs, any arithmetic or small-divisor hypothesis,
and the word \emph{resonance} appears here only as an absence: in
\cref{nab:thm:frobenius} no spectral resonance operator is inverted, because
positivity forces the exponent-zero part of the residue to vanish. Other
reports in this collection use ``resonance'' in several mutually inequivalent
senses; none of them is in force here, and no theorem of this report may be
combined with one of those hypotheses by analogy.
\end{remark}

\begin{remark}[Which exponential and which logarithm]
\label{nab:rem:explog}
The matrix $\exp$ and $\log$ in this report are formal one-variable
substitutions from \cref{nab:lem:inverse}, computed coefficientwise on a fixed
ordinary domain. They are mutually inverse maps
$\exp:\Mat_r(\JJ(V))\to\GG(V)$ and
$\log:\GG(V)\to\Mat_r(\JJ(V))$. For commuting arguments the usual
addition identities hold; they are not imposed on noncommuting arguments.
Ordinary scalar logarithms such as $\log 2$ and the branch $\Log z$ retain
their ordinary complex meanings. No matrix exponential here is the
Ehrlich--Kaplan surcomplex exponential, or any radial or quaternionic
extension of it, whose homomorphism properties are separate questions
belonging to separate reports. The scalar identity
$\Pol(G)=\exp(\pi_-\log G)$ of \eqref{nab:eq:scalarPol} is used only in rank
one, and \cref{nab:prop:logdiag} shows what happens if one tries to use it in
rank three.
\end{remark}

\subsection{Surcomplex evaluation on a halo}

At an ordinary point $z_0\in V$ and an $\varepsilon\in\mm$, a positive family
is evaluated by the ordinary Taylor series of its coefficients:
\begin{equation}\label{nab:eq:halo}
 F^\sharp(z_0+\varepsilon)
 =\sum_{\gamma}\sum_{k\geq0}
   t^\gamma\frac{F_\gamma^{(k)}(z_0)}{k!}\varepsilon^k,
\end{equation}
the exponent-zero term being allowed separately in the integral case. If $S$
is the Hahn support and $T=\supp\varepsilon$, the resulting support lies in
$S+\angles T$, and each fixed exponent receives only finitely many
contributions by \cref{nab:lem:neumann}. So our matrices do give matrices on
finite surcomplex thickenings of ordinary charts, entry by entry. This is
offered only to explain why the category deserves to be called surcomplex.
The ordinary complex base is not being replaced by a fine-topological
surcomplex manifold, and no monodromy or gluing statement below is obtained by
transport along a fine-continuous surcomplex loop.

\subsection{Ordinary paths, not fine-surreal paths}
\label{nab:sub:ordinarypaths}

For a piecewise smooth ordinary path $c$ in an ordinary domain, integrals of
coefficient one-forms and their iterated integrals are ordinary complex
integrals; their values are then assembled as Hahn coefficients. Paths, the
universal cover, iterated integrals and Mittag--Leffler are all ordinary
complex. Only the assembly of coefficients is Hahn. We assert no convergence
of partial Hahn sums in the fine topology of $\SC$, and
\cref{nab:ex:boundedsupport} shows that such partial sums can fail to converge
even in the valuation topology. This is the same domain separation stated in
the repository's differential-equations report \cite{RepoDiff}, whose wording
we adopt rather than re-argue: the loop is an ordinary complex path,
continuation is coefficientwise, and the result is a matrix of Hahn constants.

\section{Support calculus without any convergence assumption}
\label{nab:sec:support}

\begin{definition}[Strong summability]\label{nab:def:strong}
A set-indexed family of Hahn series, or of matrices of Hahn series, is
\emph{strongly summable} if the union of its supports is well ordered and each
exponent occurs in the support of only finitely many members. Its sum is then
the finite coefficientwise sum at each exponent. This is not
valuation-topology convergence, which is never used below.
\end{definition}

\begin{lemma}[Neumann support calculus; classical input]
\label{nab:lem:neumann}
Let $E\subseteq\Gamma_{>0}$ be well ordered and let
$M=\angles E=\{0\}\cup\{e_1+\cdots+e_n:n\geq1,\ e_j\in E\}$. Then $M$ is well
ordered, and for each $\gamma\in\Gamma$ only finitely many finite words in
$E$, across all word lengths, have sum $\gamma$. Consequently, if
$S\subseteq\Gamma$ is well ordered then $S+M$ is well ordered, a fixed element
has only finitely many decompositions with one summand in each of two
well-ordered sets, and every fixed coefficient of an expression formed by
adjoining finitely many positive-support factors receives only finitely many
contributions.
\end{lemma}

This is the positive-support form of Neumann's lemma
\cite{Neumann}, with the maximally complete field background in
\cite{Poonen}. We use it as an established combinatorial input and do not
rename it. It is strictly stronger than the statement that all exponents are
positive, and the finite-word clause across all word lengths is exactly what
powers noncommutative Dyson expansions and geometric inverses below. Both
source manuscripts of this report state and use it independently, in the same
form.

\begin{remark}[No hidden Archimedean or cofinality assumption]
\label{nab:rem:nonarch}
A least positive input exponent $\delta$ need not have the property that
$n\delta$ eventually exceeds a prescribed $\gamma$. In $\Z^2$ with
lexicographic order, $n(0,1)<(1,0)$ for every ordinary integer $n$. In
$\Gamma=\Q+\Q\omega$ with its induced surreal order, $n<\omega$ for every
ordinary positive integer $n$, so repeatedly increasing a valuation by $1$
never reaches $\omega$; nevertheless $\sum_{n\geq0}t^n$ is a valid Hahn sum
and the support calculus works equally well at the exponents $\omega$ and
$\omega+1$, which is where the correction of \cref{nab:prop:mixed} will sit.
No proof below infers coefficient finiteness from that false cofinality.
Every recursion is a well-founded recursion on a well-ordered support set.
\end{remark}

\begin{lemma}[Positive inversion and formal functions]\label{nab:lem:inverse}
Let $V$ be an ordinary domain and $A\in\Mat_r(\JJ(V))$, with
$\supp A\subseteq E$ for a positive well-ordered $E$. Then
\[
 (I_r+A)^{-1}=\sum_{k\geq0}(-A)^k,
 \qquad
 \exp A=\sum_{k\geq0}\frac{A^k}{k!},
 \qquad
 \log(I_r+A)=\sum_{k\geq1}\frac{(-1)^{k+1}}{k}A^k
\]
are strongly summable coefficientwise on $V$, with nonconstant supports in
$\angles E\setminus\{0\}$. The inverse is two-sided, and $\exp$ and $\log$ are
mutually inverse as one-variable formal substitutions. They do not make matrix
multiplication additive. The same statements hold with coefficients in any
associative complex algebra for which the finitely many operations in each
displayed coefficient are defined.
\end{lemma}

\begin{proof}
Expand an entry of $A^k$: its terms are indexed by a word of $k$ positive
exponents together with a path through $r$ matrix indices. At a fixed total
exponent only finitely many exponent words occur by \cref{nab:lem:neumann},
and each word admits only finitely many matrix-index paths. Hence every
displayed coefficient sum is finite and its value is holomorphic on the
original $V$, and all supports lie in $\angles E$. Ordinary formal identities
in one variable, and the geometric-series cancellation, may therefore be
checked coefficientwise, because at a fixed exponent all sufficiently long
words fail to contribute. Noncommuting independent arguments do not satisfy
the scalar exponential addition law.
\end{proof}

\begin{lemma}[Finite decomposition inside the generated monoid]
\label{nab:lem:finiteclosure}
Put $S=\angles E\setminus\{0\}$ for a positive well-ordered $E$. Each
$\gamma\in\Gamma$ has only finitely many ordered decompositions into members
of $S$, and only finitely many members of $S$ can occur as a part in such a
decomposition of a fixed $\gamma$.
\end{lemma}

\begin{proof}
Choose an $E$-word representing each member of $S$. Flattening a word in $S$
of total weight $\gamma$ gives an $E$-word of that weight, and there are
finitely many such by \cref{nab:lem:neumann}. A fixed finite word has only
finitely many partitions into nonempty consecutive blocks, and every original
$S$-word is recovered from one of these finitely many possibilities. Taking
all parts of those words gives the second assertion.
\end{proof}

\Cref{nab:lem:finiteclosure} is what makes the effectivity statement
\cref{nab:prop:dependence} possible; it is not needed for
\cref{nab:lem:inverse}.
\section{The ordinary analytic inputs}\label{nab:sec:ordinary}

Everything analytic in this report is an ordinary complex theorem applied at
\emph{one Hahn exponent at a time}. That is the reason the criteria exist at
all: the ordinary problems are always solvable, and the obstruction is that
their outputs across exponents must assemble into one legitimate Hahn
section. The two parts use two different packagings of the ordinary input, and
both are kept, because each theorem below uses its own.

\subsection{A fixed-disk Laurent splitting}

Fix $D=\{|u|<\rho\}$ and $D^*=D\setminus\{0\}$. A function $f\in\OO(D^*)$ has
the Laurent decomposition
\begin{equation}\label{nab:eq:Laurent}
 f=f_-+f_+,\qquad
 f_-=\sum_{k\geq1}a_{-k}u^{-k},\qquad
 f_+=\sum_{k\geq0}a_ku^k.
\end{equation}
Write $\PP$ for the algebra of all such negative parts and $\pi_-,\pi_+$ for
the projections, acting entrywise on matrices.

\begin{lemma}[The ordinary splitting algebras]\label{nab:lem:laurent}
There is a direct-sum decomposition of complex vector spaces
$\OO(D^*)=\PP\oplus\OO(D)$. Both summands are closed under multiplication,
although $\PP$ has no multiplicative identity; their intersection is zero; and
the decomposition does not shrink $D$.
\end{lemma}

\begin{proof}
For every $s\in(0,\rho)$, Cauchy's formula bounds the negative Laurent
coefficients by $|a_{-k}|\leq M_ss^k$, so the series in $1/u$ converges for
$|u|>s$ with $s$ arbitrarily small, and vanishes at infinity. Conversely a
function holomorphic on $\C\setminus\{0\}$ and vanishing at infinity has only
negative Laurent powers. Products retain these properties. A function in the
intersection extends across zero and across infinity, hence is constant on the
sphere with value zero at infinity. The positive part converges on the
original disk of radius $\rho$.
\end{proof}

The negative part may have an essential singularity at zero; no bound on pole
orders is imposed. The explicit obstruction families below use only simple
poles.

\subsection{Isolated singular parts: the Part~I packaging}

\begin{lemma}[Isolated singular parts on the plane]\label{nab:lem:mlplane}
Let $A\subset\C$ be closed and discrete and let $p_a$ be a prescribed negative
Laurent part in $z-a$ at each $a\in A$. There exists $f\in\OO(\C\setminus A)$
such that $f-p_a$ extends holomorphically through each disk $D_a$. One may
additionally prescribe $f(z_*)$ at one fixed $z_*\notin A$.
\end{lemma}

\begin{proof}
The finite case is a finite sum. Otherwise enumerate $A=\{a_n\}$ with
$|a_n|\to\infty$ and choose radii $R_n\to\infty$ with $R_n<|a_n|$ for all
large $n$. Each $p_{a_n}(z-a_n)$ is holomorphic on $|z|<|a_n|$; choose a
Taylor polynomial $q_n$ at zero with
$\sup_{|z|\leq R_n}|p_{a_n}(z-a_n)-q_n(z)|<2^{-n}$. After separating finitely
many early indices, $f(z)=\sum_n(p_{a_n}(z-a_n)-q_n(z))$ converges normally on
every compact subset of $\C\setminus A$. Near a fixed $a_m$ all terms except
the $m$th form a normally convergent holomorphic series, so the required
negative part is $p_{a_m}$, and the differences extend over the full $D_{a_m}$
by uniqueness. Adding a constant fixes $f(z_*)$.
\end{proof}

The proof allows essential isolated singularities, because it approximates the
entire negative part away from its centre, and uses no bound on pole orders.
The same correcting-polynomial mechanism appears in the scalar Cousin step of
the repository's global-divisor report \cite{RepoGlobal}.

\begin{lemma}[Open-surface version]\label{nab:lem:mlsurface}
The conclusion of \cref{nab:lem:mlplane} holds on any connected open Riemann
surface $X$, with $A$ closed and discrete, pairwise disjoint coordinate disks
$D_a$, and negative Laurent parts in the chosen coordinates.
\end{lemma}

\begin{proof}
An open Riemann surface is Stein by the Behnke--Stein theorem \cite{BS}. The
data $p_a$ define an ordinary additive holomorphic cocycle on the cover
$X\setminus A$, $D_a$. Cartan's theorem B gives $H^1(X,\OO)=0$ and hence a
splitting of that cocycle; in the chosen sign convention the central component
is the desired $f$. A trivial additive torsor has a section expressed in the
original charts, so no acyclicity assumption on this particular cover is
needed. A constant again prescribes one value. See \cite{Cartan,Schmitt} for
the ordinary coherent-sheaf input.
\end{proof}

On the plane, \cref{nab:lem:mlplane} makes the analytic existence step of
\cref{nab:thm:criterion} self-contained. The general open-surface statement
uses exactly the two established ordinary theorems named in its proof, and not
an unproved theorem B for a Hahn coefficient sheaf.

\subsection{The period map and a linear section: the Part~II packaging}

Let $D\subset\C$ be closed and discrete. Every compact set meets $D$ in a
finite set, so $D$ is at most countable, and $U=\C\setminus D$ is connected.

\begin{lemma}[Free based meridians]\label{nab:lem:free}
There is a collection of based positively oriented meridians
$(\ell_d)_{d\in D}$ forming a free basis of $\pi_1(U,p)$; every group element
is a finite word in these meridians and their inverses.
\end{lemma}

\begin{proof}
Exhaust $\C$ by nested disks centred at $p$ with boundaries missing $D$. Each
disk contains finitely many punctures and has free fundamental group generated
by based meridians. When the disk is enlarged, the old meridians can be
retained and meridians for the new punctures added: cutting the annular region
along arcs to the newly included punctures, or applying van Kampen, exhibits
the old group as a free factor. Choose these bases compatibly. A loop and any
specified homotopy have compact image, hence occur in a sufficiently large
disk, so the fundamental group of the union is free on the union of the
compatible bases. For $D=\varnothing$ the group is trivial.
\end{proof}

We fix such a basis throughout Part~II. The numerical matrices assigned to
based meridians depend on the chosen tails and on the base point, and that
dependence is not suppressed. So that $dY=\Omega Y$ yields a homomorphism for
matrix multiplication, $\alpha\beta$ means traversal of $\beta$ first and then
$\alpha$. All meridians are positively oriented.

Let $\LL(D)$ be the complex vector space of meromorphic one-forms on $\C$ with
at most simple poles at the points of $D$ and no other finite poles. There is
\emph{no condition at infinity}. Define
\begin{equation}\label{nab:eq:periodmap}
 P:\LL(D)\longrightarrow\C^D,
 \qquad (P\alpha)_d=\int_{\ell_d}\alpha=2\pi i\Res_d\alpha.
\end{equation}
The target is the full \emph{product} of complex vector spaces, not the direct
sum. That is the whole reason no summability across punctures is needed in
\cref{nab:thm:RH}.

\begin{lemma}[A linear period right inverse]\label{nab:lem:period}
The map $P$ is surjective, and there is a complex-linear
$\sigma:\C^D\to\LL(D)$ with $P\sigma=\id$. For finite $D$ one may take
\begin{equation}\label{nab:eq:finitesection}
 \sigma\bigl((v_d)_{d\in D}\bigr)
 =\frac1{2\pi i}\sum_{d\in D}v_d\frac{dz}{z-d}.
\end{equation}
For $D=\varnothing$ the right inverse has zero image.
\end{lemma}

\begin{proof}
Prescribing the principal part $v_d\,dz/(2\pi i(z-d))$ at each $d$ and
applying the ordinary Mittag--Leffler theorem proves surjectivity; the
elementary mechanism is the correcting-polynomial argument of
\cref{nab:lem:mlplane}, applied to the one-form data. That construction need
not be linear in the input, because the chosen polynomial degrees may depend
on it. To obtain a linear section, split the surjection of complex vector
spaces: choose a basis of $\C^D$, choose one preimage of each basis element,
and extend linearly. This uses the axiom of choice for vector spaces. For
finite $D$ the displayed finite sum is an explicit linear section.
\end{proof}

\begin{remark}[What is and is not required of the section]
\label{nab:rem:section}
For infinite $D$ we assert neither continuity nor a norm estimate nor an
effective algorithm for $\sigma$; \cref{nab:prop:discontinuous} shows that
continuity is in fact impossible for the natural topologies. In particular
$\sum_dv_d/(z-d)$ cannot simply be written without a convergence correction.
The theorem needs only an ordinary complex-linear splitting. The section is
chosen once, independently of the Hahn exponent, and is applied
\emph{entrywise} to matrix data; that last point is what makes the traceless
refinement \cref{nab:thm:SL} true.
\end{remark}

Both packagings are the ordinary Mittag--Leffler theorem, and neither is
derived from the other here. \Cref{nab:lem:mlplane} gives, for each exponent
separately, one function with prescribed negative parts of unbounded order;
\cref{nab:lem:period} gives a single \emph{linear} choice of simple-pole
one-form, valid simultaneously for all exponents, at the cost of the axiom of
choice and of continuity. Part~I needs the first because its prescribed local
data are arbitrary negative Laurent matrices; Part~II needs the second because
uniqueness of its normalized realizer is exactly injectivity of $P$ on
$\im\sigma$.

\section{Local positive factorizations at a puncture}\label{nab:sec:local}

Two local normal forms are used, one per part. They are different
normalizations of different objects and neither is a consequence of the other.

\subsection{The normalized polar factorization}

For the chosen disk and coordinate put
\[
 \mathscr P_{r,\Gamma}=I_r+\Mat_r\bigl(\PP((t^\Gamma))_{>0}\bigr),
\]
meaning that the nonconstant coefficients lie in $\PP$ and have positive well
ordered support. This is a group: nonconstant coefficients of a product or of
a positive inverse are finite sums of products of negative Laurent parts, and
each such product still vanishes at auxiliary infinity.

\begin{theorem}[Local normalized factorization]\label{nab:thm:local}
Let $G\in\GG(D^*)$ and $S=\supp(G-I_r)$. There exist unique matrices
\begin{equation}\label{nab:eq:localfactor}
 P\in\mathscr P_{r,\Gamma},\qquad H\in\GG(D),\qquad G=PH,
\end{equation}
and both nonconstant supports lie in $M=\angles S\setminus\{0\}$. We write
$P=\Pol(G)$ and $H=\Reg(G)$. Every coefficient of $H$ is holomorphic on the
\emph{original} disk, not on an exponent-dependent smaller disk.
\end{theorem}

\begin{proof}
The case $S=\varnothing$ is immediate. Otherwise write
$G=I_r+\sum_{\gamma>0}g_\gamma t^\gamma$,
$P=I_r+\sum_{\gamma\in M}p_\gamma t^\gamma$,
$H=I_r+\sum_{\gamma\in M}h_\gamma t^\gamma$, with $g_\gamma=0$ at exponents
not originally present, and use well-founded recursion on $M$. Once all
smaller coefficients are chosen, define
\begin{align}
 r_\gamma&=g_\gamma-
   \sum_{\substack{\eta+\theta=\gamma\\ \eta,\theta>0}}p_\eta h_\theta,
   \label{nab:eq:factorrec1}\\
 p_\gamma&=\pi_-(r_\gamma),\qquad h_\gamma=\pi_+(r_\gamma).
   \label{nab:eq:factorrec2}
\end{align}
Every summand involves strictly smaller exponents and the sum is finite by
\cref{nab:lem:neumann}, so each step produces a negative Laurent matrix and a
matrix holomorphic on $D$. The series have supports in $M$ and satisfy $PH=G$
coefficientwise; their inverses exist by \cref{nab:lem:inverse}.

For uniqueness, suppose $P_1H_1=P_2H_2$, so $P_2^{-1}P_1=H_2H_1^{-1}$. The
nonconstant coefficients on the left lie in $\PP$ and those on the right are
holomorphic on $D$. After placing the finitely many input supports in a common
positive generated monoid, coefficientwise comparison and
\cref{nab:lem:laurent} force both sides to equal $I_r$.
\end{proof}

\begin{proposition}[Right-gauge invariance]\label{nab:prop:rightgauge}
If $Q\in\GG(D)$ then
\begin{equation}\label{nab:eq:rightgauge}
 \Pol(GQ)=\Pol(G).
\end{equation}
The analogous statement for multiplication on the left is generally false.
\end{proposition}

\begin{proof}
The factorization $GQ=\Pol(G)(\Reg(G)Q)$ has the defining properties of
\cref{nab:thm:local}, and uniqueness gives the equality. For failure on the
left, choose $\delta>0$, $s=t^\delta$, $G=I_3+(s/u)E_{23}$ and
$Q=I_3+sE_{12}$. Direct multiplication gives
\[
 QG=\bigl(I_3+(s/u)E_{23}+(s^2/u)E_{13}\bigr)Q.
\]
Here $E_{ij}$ are matrix units. The first factor is polar and $Q$ is regular,
so it is $\Pol(QG)$, which
differs from $\Pol(G)=G$. This exhibits the multiplication order explicitly.
\end{proof}

\begin{corollary}[Unitriangular factors]\label{nab:cor:unitriangular}
If $G$ is upper unitriangular then so are $\Pol(G)$ and $\Reg(G)$. If a fixed
pattern of strictly upper-triangular entries is closed under matrix
multiplication, the factors preserve that pattern too.
\end{corollary}

\begin{proof}
Both Laurent projections preserve zero entries, and every nonlinear term in
\eqref{nab:eq:factorrec1} is a product of previously constructed matrices in
the specified multiplicatively closed pattern. Induction finishes the proof.
\end{proof}

\subsection{Where the obstruction is manufactured}

For a single positive parameter, writing
$G=I_r+tg_1+t^2g_2+t^3g_3+\cdots$, the recursion begins
\begin{align*}
 p_1&=\pi_-g_1,& h_1&=\pi_+g_1,\\
 p_2&=\pi_-(g_2-p_1h_1),& h_2&=\pi_+(g_2-p_1h_1),\\
 p_3&=\pi_-(g_3-p_1h_2-p_2h_1).&&
\end{align*}
The second-order correction already contains the product of a polar part and a
regular part, and it can be nonzero even when $g_2=0$. With different positive
exponents at different punctures, the exponents of such products can form a
descending sequence absent from the raw singular data. That single mechanism is
the whole content of \cref{nab:thm:hidden} and \cref{nab:thm:depth}.

In rank one the commuting logarithmic parts give the closed form
\begin{equation}\label{nab:eq:scalarPol}
 \Pol(G)=\exp\bigl(\pi_-\log G\bigr),
 \qquad
 \Reg(G)=\exp\bigl(\pi_+\log G\bigr),
\end{equation}
with the projections acting coefficientwise and the scalar exponential
addition law verifying the factorization. The same argument applies to a
matrix input if its two projected logarithmic parts commute. For general
matrix inputs,
replacing \cref{nab:thm:local} by \eqref{nab:eq:scalarPol} is invalid; the
algebraic Birkhoff literature describes the associated
Baker--Campbell--Hausdorff corrections in a filtered setting
\cite{EFGK,EFMP}, and our direct recursion avoids expanding that formula.
\Cref{nab:prop:logdiag} exhibits the exact numerical discrepancy in rank
three.

\subsection{The positive Frobenius form}

The next theorem is the local normal form used in Part~II. It explains why
residues still control \emph{local} monodromy while being unreadable directly
from the prescribed \emph{based} monodromy matrices. It is a
support-controlled version of the elementary nonresonant regular-singular
construction, not a new claim about classical Frobenius theory.

\begin{theorem}[A common-disk positive Frobenius factorization]
\label{nab:thm:frobenius}
Let $\Delta$ be an ordinary disk about $0$ and suppose a positive Hahn
connection on $\Delta\setminus\{0\}$ has the form
\begin{equation}\label{nab:eq:localconnection}
 \Omega=R\frac{dz}{z}+B(z)\,dz,
\end{equation}
with $R\in\Mat_r(\mm)$ and $B\in\Mat_r(\JJ(\Delta))$ sharing one positive well
ordered support $E$. Then there is a unique
$H\in I_r+\Mat_r(\JJ(\Delta))$ with $H(0)=I_r$ such that, on a branch of the
logarithm,
\begin{equation}\label{nab:eq:frobenius}
 Y(z)=H(z)\exp(R\Log z)
\end{equation}
solves $dY=\Omega Y$. The supports of $H-I_r$ and $H^{-1}-I_r$ lie in
$\angles E\setminus\{0\}$, and the local monodromy based at
$z_0\in\Delta\setminus\{0\}$ is
\begin{equation}\label{nab:eq:localmonodromy}
 H(z_0)\exp(2\pi iR)H(z_0)^{-1}.
\end{equation}
\end{theorem}

\begin{proof}
The required equation is $H'=[R,H]/z+BH$ with $H(0)=I_r$. Construct its
coefficients recursively on $\angles E\setminus\{0\}$: at weight $\gamma$ one
must solve
\[
 H_\gamma'=B_\gamma+
 \sum_{\substack{\alpha+\beta=\gamma\\ \alpha,\beta>0}}
  \Bigl(B_\alpha H_\beta+\frac{[R_\alpha,H_\beta]}{z}\Bigr),
 \qquad H_\gamma(0)=0,
\]
a finite sum. The previously constructed $H_\beta$ vanish at $0$, so
$H_\beta/z$ is holomorphic on the entire disk; the right-hand side is
therefore holomorphic on that same disk, and integration from $0$ defines a
unique holomorphic $H_\gamma$ vanishing there. The resulting family has the
claimed support and its inverse is given by \cref{nab:lem:inverse}. A
least-difference argument proves uniqueness.

Because $R$ has positive support, $\exp(R\Log z)$ is a strongly summable
family of ordinary holomorphic functions on any logarithm branch, and
coefficientwise differentiation gives
$(\exp(R\Log z))'=R\exp(R\Log z)/z$; the equation for $H$ then verifies
\eqref{nab:eq:frobenius}. Continuing once positively around $0$ increases
$\Log z$ by $2\pi i$ and multiplies the fundamental matrix on the right by
$\exp(2\pi iR)$. Normalizing at $z_0$ conjugates that factor by
$H(z_0)\exp(R\Log z_0)$, and the exponential factor commutes with
$\exp(2\pi iR)$ and cancels, giving \eqref{nab:eq:localmonodromy}.
\end{proof}

There is no division by a spectral resonance operator in this proof: the
exponent-zero coefficient of the residue is zero, so the equation at each new
exponent is an ordinary primitive problem. Positivity is doing real work here
(compare \cref{nab:rem:resonance}).

If a meridian is based at a distant point $p$, its tail contributes a further
transport conjugation to \eqref{nab:eq:localmonodromy}, and with several
punctures these conjugations interact. Choosing
$R_d=(2\pi i)^{-1}\log M_d$ separately at each puncture therefore need not
solve the based inverse problem. The explicit failure appears as a mixed
commutator term in \cref{nab:prop:mixed}.
\part*{Part I. Multiplicative gluing: the polar-support criterion}
\addcontentsline{toc}{section}{\textbf{Part I. Multiplicative gluing: the polar-support criterion}}

\section{The sharp global matrix Cousin criterion}\label{nab:sec:cousin}

Let $X$ be a connected open Riemann surface, $A\subset X$ closed and discrete,
and choose pairwise disjoint coordinate disks $(D_a,u_a)$ with $u_a(a)=0$. Put
$U_0=X\setminus A$ and $W_a=D_a\setminus\{a\}$; the punctured surface $U_0$ is
connected. For each $a$, \emph{independently}, let $G_a\in\GG(W_a)$. The disks
do not meet one another, so this is a complete positive matrix cocycle on the
displayed cover, with the inverse transition on the reversed overlap. There
need not be one well-ordered set containing the supports of all the $G_a$.

\begin{definition}[Normalized polar support]\label{nab:def:polsupport}
Factor $G_a=P_aH_a$ by \cref{nab:thm:local} and define
\begin{equation}\label{nab:eq:polsupport}
 \EU(G)=\bigcup_{a\in A}\supp(P_a-I_r)\subseteq\Gamma_{>0}.
\end{equation}
The datum is \emph{polar-admissible} when $\EU(G)$ is well ordered. No sum
over $a$ is intended in \eqref{nab:eq:polsupport}: it is a set-theoretic union
of supports, taken across all matrix entries.
\end{definition}

\begin{theorem}[Exact nonabelian Cousin criterion: Criterion P]
\label{nab:thm:criterion}
The following are equivalent.
\begin{enumerate}[label=\textup{(\roman*)}]
\item There exist $B_0\in\GG(U_0)$ and $B_a\in\GG(D_a)$ with
$G_a=B_0B_a^{-1}$ on every $W_a$.
\item The set $\EU(G)$ is well ordered.
\end{enumerate}
When these hold, let $M=\angles{\EU(G)}$. One can choose
\begin{equation}\label{nab:eq:centralsupport}
 \supp(B_0-I_r)\subseteq M\setminus\{0\},
\end{equation}
and in the local factorizations $B_0=P_aQ_a$ the matrices $Q_a-I_r$ can also
be supported in $M\setminus\{0\}$. The disk trivializers are
$B_a=H_a^{-1}Q_a$ and hence have individually legitimate supports. A value
$B_0(z_*)=I_r$ may be imposed at any chosen $z_*\in U_0$.
\end{theorem}

\begin{proof}
\emph{Necessity.} Suppose $G_a=B_0B_a^{-1}$. By right-gauge invariance
\eqref{nab:eq:rightgauge},
\[
 P_a=\Pol(G_a)=\Pol\bigl(B_0|_{W_a}\bigr).
\]
Let $S=\supp(B_0-I_r)$. The local support estimate in \cref{nab:thm:local} is
independent of the puncture and of the Laurent coefficients of the input, so
$\supp(P_a-I_r)\subseteq\angles S\setminus\{0\}$ for every $a$. That one
generated monoid is well ordered by \cref{nab:lem:neumann}, and $\EU(G)$ is a
subset of it.

\emph{Sufficiency.} Assume $E=\EU(G)$ is well ordered and put $M=\angles E$.
We construct a global $F\in\GG(U_0)$ and local $Q_a\in\GG(D_a)$ with
\begin{equation}\label{nab:eq:Fgoal}
 F=P_aQ_a
\end{equation}
and every nonconstant coefficient supported in $M$. Write
$P_a=I_r+\sum_{\gamma>0}p_{a,\gamma}t^\gamma$,
$F=I_r+\sum_{\gamma>0}f_\gamma t^\gamma$,
$Q_a=I_r+\sum_{\gamma>0}q_{a,\gamma}t^\gamma$. At a positive $\gamma\in M$,
with all smaller coefficients constructed, define on $W_a$
\begin{equation}\label{nab:eq:globalrec}
 R_{a,\gamma}=p_{a,\gamma}+
 \sum_{\substack{\eta+\theta=\gamma\\ \eta,\theta>0}}p_{a,\eta}q_{a,\theta},
\end{equation}
a finite sum. Apply \cref{nab:lem:mlsurface} entry by entry to the prescribed
ordinary negative Laurent matrices $\pi_-R_{a,\gamma}$. It supplies a single
matrix $f_\gamma\in\Mat_r(\OO(U_0))$ such that $f_\gamma-R_{a,\gamma}$ extends
holomorphically across $D_a$ for every $a$; set
\begin{equation}\label{nab:eq:qrec}
 q_{a,\gamma}=f_\gamma-R_{a,\gamma}.
\end{equation}
The prescribed zero value at $z_*$ can be imposed on every $f_\gamma$.

Recursion on the well-ordered set $M\setminus\{0\}$ constructs all
coefficients. The scalar analytic step changes coefficient functions but
introduces no new Hahn exponents, and every coefficient of every $Q_a$ is
holomorphic on its original disk. The resulting Hahn matrices satisfy
\eqref{nab:eq:Fgoal} coefficientwise and are invertible by
\cref{nab:lem:inverse}. Finally $G_a=P_aH_a$ and $F=P_aQ_a$ give
$G_a=F(H_a^{-1}Q_a)^{-1}$; take $B_0=F$ and $B_a=H_a^{-1}Q_a$. The support of
$H_a-I_r$ lies in the monoid generated by the original support of $G_a-I_r$
and the support of $Q_a-I_r$ lies in $M$; a finite union of two well-ordered
positive sets generates a legitimate monoid, so $B_a$ is a well-defined
section on $D_a$. There is no requirement that the disk supports have a
globally well-ordered union.
\end{proof}

\subsection{Why the converse is sharp}

An infinite puncture set does not by itself prevent a solution: at each fixed
exponent the ordinary analytic problem is always solvable, and the obstruction
is not a divergent ordinary sum, since the ordinary coefficient construction
uses normal convergence independent of the formal exponents. The failure occurs
exactly when the \emph{necessary normalized polar information} cannot fit into
a single well-ordered support.

Write $E_{\mathrm{raw}}=\bigcup_a\supp(G_a-I_r)$. If this set is well
ordered, \cref{nab:thm:local} puts every $P_a-I_r$ in
$\angles{E_{\mathrm{raw}}}\setminus\{0\}$, so Criterion~P gives a splitting.
Thus well-ordering of the full raw union is a sufficient condition. It is
not necessary: $G_a$ may itself be holomorphic on $D_a$, in which case
$P_a=I_r$ and the solution
$B_0=I_r$, $B_a=G_a^{-1}$ exists even if those raw supports have a descending
union. In contrast, the raw singular-coefficient union
$\bigcup_a\supp\bigl(\pi_-(G_a-I_r)\bigr)$ fails both directions:
discarding raw holomorphic coefficients can lose polar information created
by their products with retained terms, and retained singular terms can cancel
in the normalization. \Cref{nab:sec:rankthree} gives both examples.

\subsection{The space of solutions}

\begin{proposition}[Right torsor of splittings]\label{nab:prop:torsor}
For a polar-admissible datum, all splittings are obtained from any one by
\[
 (B_0,B_a)\longmapsto(B_0T,\,B_aT|_{D_a}),\qquad T\in\GG(X),
\]
and the action is free and transitive. Under the normalization $B_0(z_*)=I_r$,
use the subgroup with $T(z_*)=I_r$.
\end{proposition}

\begin{proof}
Right multiplication by a common $T$ preserves $B_0B_a^{-1}$. Conversely two
splittings give $B_0^{-1}B_0'=B_a^{-1}B_a'$ on $W_a$; the central quotient
extends through every disk by that equality, and the sheaf observation of
\cref{nab:sub:conventions} yields a global positive section $T$ with one
legitimate support on connected $X$. It is the unique element relating the two
splittings.
\end{proof}

This describes all solutions once a splitting exists. It is not a
classification of the nonsplit objects; see \cref{nab:sec:scope}.

\section{Invariance, and the scalar case that is recovered}
\label{nab:sec:invariance}

\begin{corollary}[Removal of arbitrarily supported regular data]
\label{nab:cor:localgauges}
For independently chosen $R_a\in\GG(D_a)$, the families $(G_a)$ and
$(G_aR_a)$ have the same normalized polar factors and the same solvability
status. No well-ordering condition on $\bigcup_a\supp(R_a-I_r)$ is needed.
\end{corollary}

\begin{proof}
Immediate from \cref{nab:prop:rightgauge} and \cref{nab:thm:criterion}.
\end{proof}

This is stronger than invariance under finitely many transformations: the
infinitely many regular parts are genuinely irrelevant, but only \emph{after}
they have been removed in the correct multiplication order.

\begin{corollary}[Coordinate independence of admissibility]
\label{nab:cor:coordinates}
Changing the local coordinates, or shrinking the disks, can change individual
normalized polar coefficients but cannot change whether their global support
union is well ordered.
\end{corollary}

\begin{proof}
Existence of the splitting in \cref{nab:thm:criterion} is invariant under
these changes. Shrinking a disk cannot turn a previously nonextendible germ
into an extendible one: a local splitting on smaller disks yields extension of
the same central matrix, hence a splitting on the original disks by analytic
continuation of $G_a^{-1}B_0$. Alternatively, re-expressing a previously
constructed polar factor in the new coordinate and applying
\cref{nab:thm:local} puts every new support in the monoid generated by the old
support union, and the inverse coordinate change gives the converse.
\end{proof}

\begin{corollary}[Absoluteness under ordered extension]
\label{nab:cor:extension}
Let $\Gamma\hookrightarrow\Delta$ be an order-preserving embedding of
set-sized ordered abelian groups. A datum defined over $\Gamma$ splits in the
positive matrix category over $\Delta$ if and only if it splits over $\Gamma$.
\end{corollary}

\begin{proof}
The recursive local polar factor computed in $\Gamma$ is also a factor in
$\Delta$ and remains the unique normalized factor, so its support union is the
image of $\EU(G)$. An order embedding preserves and reflects well-ordering of
a subset. Apply \cref{nab:thm:criterion} in both groups.
\end{proof}

In particular the obstructions below cannot be removed by allowing more
surreal exponents in the same integral positive-matrix problem. This is not a
claim about passing to another analytic category, and in particular not a
claim about admitting negative-valuation gauges; see
\cref{nab:sec:bundles} and item~\ref{nab:nc:localization} of
\cref{nab:sec:scope}.

\begin{proposition}[Recovery of the rank-one singular-support test]
\label{nab:prop:rankone}
Let $r=1$, put $L_a=\log G_a$, and define the logarithmic singular support
\[
 \Sigma(G)=\bigcup_a\supp(\pi_-L_a).
\]
Then
\begin{equation}\label{nab:eq:scalarequiv}
 \EU(G)\ \text{well ordered}
 \quad\Longleftrightarrow\quad
 \Sigma(G)\ \text{well ordered}.
\end{equation}
\end{proposition}

\begin{proof}
By \eqref{nab:eq:scalarPol}, $\log P_a=\pi_-L_a$ is supported in the positive
monoid generated by $\EU(G)$, which gives the forward direction; conversely
exponentiation puts every $P_a-I_1$ inside the monoid generated by
$\Sigma(G)$. Apply \cref{nab:lem:neumann}.
\end{proof}

Combined with \cref{nab:thm:criterion}, \eqref{nab:eq:scalarequiv} says that
in rank one our criterion is exactly the exact Cousin criterion already proved
in the repository's global-divisor report --- its Theorem~13.1, stated there in
terms of the singular support of the additive data on a puncture-and-disk
cover \cite{RepoGlobal}. We cite that theorem as the case we recover; it is not
reproved here, and no priority for the scalar statement is claimed. What is
new is that the correct matrix replacement for $\pi_-\log$ is the normalized
polar factor and not the Laurent projection of a matrix logarithm: the scalar
formula is not valid for general matrix inputs because logarithms need not
turn products into sums (\cref{nab:rem:explog}), and
\cref{nab:prop:logdiag} measures the error.

\section{A rank-three obstruction invisible in the raw poles}
\label{nab:sec:rankthree}

This section contains the example that forces the distinction of
\cref{nab:sec:twocriteria}. Read it before generalizing either criterion.

\subsection{The closed factorization in the Heisenberg group}

\begin{proposition}[Closed Heisenberg factorization]
\label{nab:prop:heisenberg}
Let
\[
 G=\begin{pmatrix}1&a&b\\0&1&c\\0&0&1\end{pmatrix},
 \qquad a,b,c\in\JJ(D^*),
\]
and write $a=a_-+a_+$, $c=c_-+c_+$ by coefficientwise Laurent projection. Then
\begin{equation}\label{nab:eq:HeisenbergP}
 \Pol(G)=
 \begin{pmatrix}
  1&a_-&\pi_-(b-a_-c_+)\\ 0&1&c_-\\ 0&0&1
 \end{pmatrix},
 \qquad
 \Reg(G)=
 \begin{pmatrix}
  1&a_+&\pi_+(b-a_-c_+)\\ 0&1&c_+\\ 0&0&1
 \end{pmatrix}.
\end{equation}
\end{proposition}

\begin{proof}
All products are legitimate, being finite expressions in finitely many Hahn
series. Direct multiplication verifies $\Pol(G)\Reg(G)=G$, the left factor is
polar and the right factor is holomorphic on the full disk, so uniqueness in
\cref{nab:thm:local} identifies them as the canonical factors.
\end{proof}

The subtraction of $a_-c_+$ is the essential term: a product of a \emph{polar}
coefficient of one entry with a \emph{holomorphic} coefficient of another. Its
exponent is a sum of two input exponents that need not appear in the input.

\subsection{The obstructed datum}

Take $X=\C$, $A=\{n:n\geq2\}$, $D_n=\{|z-n|<1/4\}$, $\Gamma=\Q$, and
$u_n=z-n$. Consider
\begin{equation}\label{nab:eq:badG}
 G_n=I_3+\frac{t}{u_n}E_{12}+t^{1/n}E_{23},
\end{equation}
with $E_{ij}$ the matrix units. Each $G_n$ is a finite Hahn expression, has
determinant exactly one, and reduces to $I_3$. By
\eqref{nab:eq:HeisenbergP},
\begin{equation}\label{nab:eq:badPH}
 P_n=I_3+\frac{t}{u_n}E_{12}-\frac{t^{1+1/n}}{u_n}E_{13},
 \qquad
 H_n=I_3+t^{1/n}E_{23},
\end{equation}
and the product $P_nH_n$ has zero $(1,3)$ entry because its two contributions
cancel.

\begin{theorem}[Hidden rank-three obstruction]\label{nab:thm:hidden}
The matrices \eqref{nab:eq:badG} admit no splitting by general positive
$3\times3$ matrices. Their abelianized upper-unitriangular cocycle does split.
Their raw nonremovable coefficient support is the singleton $\{1\}$, and their
determinant cocycle is identically one.
\end{theorem}

\begin{proof}
The normalized polar support contains $\{1+1/n:n\geq2\}$, a strictly
decreasing infinite sequence, hence is not well ordered, and
\cref{nab:thm:criterion} rules out a splitting even in the full positive
matrix group.

The abelianization of the upper-unitriangular $3\times3$ group remembers the
$(1,2)$ and $(2,3)$ entries and adds them componentwise. Let
$\phi\in\OO(U_0)$ have principal part $1/(z-n)$ at every $n\geq2$ and no
singularities outside $A$; such $\phi$ exists by \cref{nab:lem:mlplane}, one
explicit choice being
\begin{equation}\label{nab:eq:phi}
 \phi(z)=\sum_{n=2}^\infty\Bigl(\frac1{z-n}+\frac1n\Bigr),
\end{equation}
whose $n$th summand is $O_K(n^{-2})$ on each compact set away from $A$. The
$(1,2)$ component splits using the central function $t\phi$ and the disk
functions $t(\phi-1/u_n)$; the $(2,3)$ component splits using central function
$0$ and disk functions $-t^{1/n}$. Finally the only original singular
coefficient is $E_{12}/u_n$ at exponent $1$, and unitriangular matrices have
determinant one.
\end{proof}

The example is stronger than a determinant obstruction, and stronger even than
the statement that a noncommutative cocycle has nontrivial abelianization: the
abelianized problem \emph{is} solvable, and the obstruction appears only when
its solutions are required to fit together in the original multiplication law.
The obstruction also lives in a matrix entry that is identically zero in the
input.

\subsection{Adding a singular term repairs the problem}

\begin{example}[The repair, which worsens the raw support]
\label{nab:ex:repair}
Define
\begin{equation}\label{nab:eq:goodG}
 \widetilde G_n=I_3+\frac{t}{u_n}E_{12}+t^{1/n}E_{23}
  +\frac{t^{1+1/n}}{u_n}E_{13}.
\end{equation}
Then
\[
 \widetilde G_n=\Bigl(I_3+\frac{t}{u_n}E_{12}\Bigr)
                \bigl(I_3+t^{1/n}E_{23}\bigr),
\]
so $\Pol(\widetilde G_n)=I_3+tE_{12}/u_n$ and
$\EU(\widetilde G)=\{1\}$, and a global splitting exists. Explicitly, with
$\phi$ as in \eqref{nab:eq:phi}, put $F=I_3+t\phi E_{12}$,
\[
 Q_n=I_3+t\Bigl(\phi-\frac1{u_n}\Bigr)E_{12},
 \qquad
 \widetilde H_n=I_3+t^{1/n}E_{23};
\]
then $F=\Pol(\widetilde G_n)Q_n$ and
$\widetilde G_n=F(\widetilde H_n^{-1}Q_n)^{-1}$, all disk factors being
holomorphic on $D_n$. Yet the raw singular support of \eqref{nab:eq:goodG}
\emph{contains} the descending set $\{1+1/n\}$.
\end{example}

So \eqref{nab:eq:badG} has raw singular support $\{1\}$ and is obstructed,
while \eqref{nab:eq:goodG} has a descending raw singular support and is
solvable. Raw singular support is therefore neither a sufficient nor a
necessary test, in the strongest possible sense: the two directions fail on
data that differ by one term.

\subsection{Why a projected matrix logarithm is not a repair}

\begin{proposition}[Logarithm diagnostic]\label{nab:prop:logdiag}
The matrix logarithm of \eqref{nab:eq:badG} is finite, namely
\[
 \log G_n=\frac{t}{u_n}E_{12}+t^{1/n}E_{23}
  -\frac12\frac{t^{1+1/n}}{u_n}E_{13},
\]
so exponentiating its negative Laurent part alone gives the $(1,3)$
coefficient $-\tfrac12t^{1+1/n}/u_n$, whereas the correct polar coefficient in
\eqref{nab:eq:badPH} is $-t^{1+1/n}/u_n$.
\end{proposition}

\begin{proof}
$G_n-I_3$ is nilpotent of order three, so the series of
\cref{nab:lem:inverse} terminates: with $a=t/u_n$ and $c=t^{1/n}$ one has
$(G_n-I_3)^2=ac\,E_{13}$ and $(G_n-I_3)^3=0$, giving
$\log G_n=(G_n-I_3)-\tfrac12ac\,E_{13}$, which is the display. Comparison with
\eqref{nab:eq:badPH} gives the factor $\tfrac12$. The symbolic check is
recorded in the companion code (\cref{nab:sec:verification}).
\end{proof}

Thus even the true matrix logarithm is not a substitute for the factorization:
one still needs the noncommutative correction, and the global criterion is
built from the actual normalized factor rather than from an \emph{ad hoc}
linearization. Note also that the discrepancy is only a factor, not a change
of exponent, so a reader who used the logarithm would find the same descending
set here --- and \cref{nab:ex:repair} shows that agreement is an accident of
this example, not a method.

\section{Obstructions at arbitrary nilpotent depth}\label{nab:sec:depth}

\subsection{A chain identity}

\begin{lemma}[Chain identity]\label{nab:lem:chain}
Let $r\geq3$, let $a$ be a negative Laurent Hahn function with positive
support, and let $c_2,\ldots,c_{r-1}$ be positive Hahn constants. Put
\[
 C=\sum_{j=2}^{r-1}c_jE_{j,j+1},\qquad H=I_r+C,\qquad G=I_r+C+aE_{12}.
\]
Since $C^{r-1}=0$,
\begin{equation}\label{nab:eq:chainP}
 P=GH^{-1}=I_r+aE_{12}(I_r+C)^{-1}
  =I_r+\sum_{k=2}^{r}(-1)^{k-2}a\Bigl(\prod_{j=2}^{k-1}c_j\Bigr)E_{1k},
\end{equation}
the empty product for $k=2$ being $1$, and \eqref{nab:eq:chainP} is the
normalized polar factorization $P=\Pol(G)$, $H=\Reg(G)$.
\end{lemma}

\begin{proof}
The finite geometric inverse of $I_r+C$ gives the displayed signs. Every
nonconstant coefficient of $P$ is a negative Laurent part, because each is a
constant multiple of $a$, while $H$ is holomorphic on the disk. Uniqueness in
\cref{nab:thm:local} identifies the factors.
\end{proof}

\subsection{The exponent design}

For $r\geq4$ and $n\geq2$ set
\begin{equation}\label{nab:eq:deepdata}
 a_n=\frac{t}{z-n},\qquad
 c_{2,n}=t^{1-1/n},\qquad
 c_{j,n}=t\ (3\leq j\leq r-2),\qquad
 c_{r-1,n}=t^{2/n},
\end{equation}
the middle range being empty for $r=4$. Let $C_n$ and $G_n$ be the matrices of
\cref{nab:lem:chain} built from these data. Every $G_n$ is finite, upper
unitriangular, and has exactly one raw singular matrix coefficient, namely
$E_{12}/(z-n)$ at exponent $1$. For $3\leq k\leq r-1$ the exponent in
$(P_n)_{1k}$ is
\begin{equation}\label{nab:eq:lowerexponents}
 \lambda_{k,n}=k-1-\frac1n,
\end{equation}
whereas the exponent in the final entry is
\begin{equation}\label{nab:eq:lastexponent}
 \lambda_{r,n}=r-2+\frac1n.
\end{equation}
For each fixed $k<r$ the exponents \eqref{nab:eq:lowerexponents} increase with
$n$ and their set has order type $\omega$; there are finitely many such sets,
so their union with $\{1\}$ is well ordered. The final exponents
\eqref{nab:eq:lastexponent} strictly decrease.

\begin{theorem}[Arbitrarily deep central support obstructions]
\label{nab:thm:depth}
For each integer $r\geq3$ there is a positive upper-unitriangular cocycle on
the puncture-and-disk cover of $\C$ with $A=\{2,3,\ldots\}$, over
$\Gamma=\Q$, such that:
\begin{enumerate}[label=\textup{(\roman*)}]
\item its original singular coefficient support is $\{1\}$ and its determinant
is one;
\item its image in the quotient of the upper-unitriangular group by
$Z_r=\{I_r+cE_{1r}:c\in\JJ\}$ is a coboundary;
\item it is not a coboundary in the full group of positive $r\times r$
matrices;
\item it remains nonsplit after every ordered value-group extension.
\end{enumerate}
For $r=3$ use \eqref{nab:eq:badG}; for $r\geq4$ use \eqref{nab:eq:deepdata}
with \cref{nab:lem:chain}.
\end{theorem}

\begin{proof}
Property (i) was checked directly above. \Cref{nab:lem:chain} and
\eqref{nab:eq:lastexponent} show that the normalized polar support contains a
strictly descending infinite sequence, so \cref{nab:thm:criterion} proves
(iii) and \cref{nab:cor:extension} proves (iv).

For (ii), let $P_n'$ be obtained from $P_n$ by deleting its $(1,r)$ entry.
Their global support union is well ordered by
\eqref{nab:eq:lowerexponents}, the case $r=3$ leaving only exponent $1$. Every
$P_n'-I_r$ belongs to the row-one algebra
\[
 V=\operatorname{span}\{E_{12},\ldots,E_{1,r-1}\},\qquad V^2=0,
\]
so the gluing problem for the $P_n'$ is additive and splits by the scalar
construction with this one well-ordered support union: solve the ordinary
principal-part problem for each coefficient of each allowed row entry, using
\cref{nab:lem:mlplane}. This gives $F'$ on $U_0$ and $Q_n'$ on $D_n$, valued
in $I_r+V$, with $F'=P_n'Q_n'$. In the quotient by $Z_r$ one has
$\overline P_n=\overline P_n'$, and since $G_n=P_nH_n$ with $H_n$ holomorphic,
\[
 \overline G_n=\overline F'\bigl(\overline H_n^{-1}\overline Q_n'\bigr)^{-1}
\]
provides the required coboundary. All factors have individually legitimate
positive supports. No matrix representative for the quotient is silently
assumed: the bars denote the ordinary group quotient by the central subgroup.
\end{proof}

Here $Z_r$ is the central subgroup with $c\in\JJ$. Over $\Gamma=\Q$ one has
$\JJ^m=\JJ$ for every positive integer $m$, since a nonzero section of least
exponent $\delta>0$ factors into $m-1$ monomials $t^{\delta/m}$ and one
remaining positive section. The elementary commutator identity
$[I+aE_{ij},I+bE_{jk}]=I+abE_{ik}$ then shows that $Z_r$ is the final nonzero
term of the lower central series of the positive upper-unitriangular group, and
every further lower-central quotient factors through the quotient in (ii). So
all those lower-depth tests pass, and the full obstruction first appears at
depth $r-1$, for arbitrarily large finite $r$.

\begin{remark}[What the depth assertion means]\label{nab:rem:depth}
The theorem does not say that a scalar equation cannot detect the final
obstruction once a noncommutative normal form has been found --- the final
coefficient is of course scalar. It says that discarding the last central
level \emph{before} normalizing loses an obstruction created by the
multiplication law. The new information sits in a genuinely higher
noncommutative compatibility condition.
\end{remark}

\section{Integral vector bundles and local invisibility}
\label{nab:sec:bundles}

\subsection{From positive matrices to actual vector bundles}

The transition matrices $G_a$ define a locally free rank-$r$ sheaf
$\mathcal E_G$ over $\II$ on the ordinary surface $X$. Its ordinary reduction
is canonically the trivial rank-$r$ holomorphic bundle, because every
transition reduces to $I_r$, and a positive splitting is a trivialization
compatible with that reduced framing.

\begin{proposition}[For triviality, the reduced framing costs nothing]
\label{nab:prop:integraltrivial}
For a bundle defined by positive transition matrices the following are
equivalent, on an arbitrary ordinary complex base and not only on a star
cover:
\begin{enumerate}[label=\textup{(\roman*)}]
\item $\mathcal E_G$ is trivial as an $\II$-module;
\item its transition cocycle splits by positive matrices.
\end{enumerate}
\end{proposition}

\begin{proof}
A positive splitting gives a trivialization. Conversely, suppose
$A_i\in\GL_r(\II(U_i))$ satisfy $G_{ij}=A_iA_j^{-1}$ and reduce modulo $\JJ$.
Since $\red G_{ij}=I_r$, the ordinary invertible matrices $\red A_i$ agree on
overlaps and glue to $A\in\GL_r(\OO(X))$. Right multiplication by the ordinary
matrix $A^{-1}$ produces $B_i=A_iA^{-1}$ with reduction $I_r$, hence
$B_i\in\GG(U_i)$ and $G_{ij}=B_iB_j^{-1}$.
\end{proof}

\begin{corollary}[Nontrivial determinant-one integral bundles]
\label{nab:cor:nontrivialbundles}
The examples of \cref{nab:thm:hidden,nab:thm:depth} define nontrivial locally
free $\II$-modules even after forgetting the upper-triangular flag and the
chosen reduced framing. Their ordinary reductions and determinant line bundles
are trivial, the transition determinants being exactly one and not merely one
modulo positive exponents.
\end{corollary}

\begin{proof}
A general integral trivialization would give a positive matrix splitting by
\cref{nab:prop:integraltrivial}, contradicting \cref{nab:thm:criterion}.
\end{proof}

\subsection{Every bounded restriction is trivial}

\begin{corollary}[Exhaustion-invisible nontriviality]\label{nab:cor:bounded}
Each integral bundle in \cref{nab:cor:nontrivialbundles} restricts trivially
to every bounded ordinary domain in $\C$. In particular it is trivial on every
disk of an increasing exhaustion of the plane, but not on the plane itself.
\end{corollary}

\begin{proof}
A bounded domain meets only finitely many of the fixed disks $D_n$, so on the
restricted cover there are only finitely many nonidentity transition matrices,
each with legitimate positive support, and their union is well ordered. Every
ordinary plane domain is Stein, so \cref{nab:thm:Stein} applies. Equivalently,
apply the finite-puncture construction directly and absorb transitions on
pieces containing no puncture into holomorphic local gauges.
\end{proof}

The failure is compatibility of infinitely many local trivializations, not
existence on any one bounded domain. Consequently a numerical experiment
involving finitely many punctures cannot establish global triviality: it will
necessarily miss these examples. This is an explicit limit on what any
implementation can decide, and it is recorded again as
item~\ref{nab:nc:blackbox} of \cref{nab:sec:scope}.

\subsection{The localization boundary}

Let $\HH=\OO((t^\Gamma))$ be the sheaf in which arbitrary monomial shifts are
allowed; the natural map $\II\to\HH$ extends scalars from $\mathcal E_G$ to a
rank-$r$ $\HH$-module. We have \emph{not} proved that the resulting module is
nontrivial, and we decline the claim rather than gesture at it. The reduction
argument of \cref{nab:prop:integraltrivial} cannot be repeated over $\HH$:
there is no reduction ring homomorphism $\HH\to\OO$ killing positive
monomials, because those monomials are units in $\HH$.

Accordingly the precise conclusions of Part~I are nonexistence of a positive
matrix splitting with common-domain Hahn entries and nontriviality of the
corresponding \emph{integral} locally free module. Whether negative-valuation changes of
frame can trivialize particular higher-rank examples is a separate question.
The scalar report proves more in rank one by a special leading-unit
normalization \cite{RepoGlobal}; no matrix analogue of that argument is
assumed here.

This distinction matters for the surcomplex interpretation. Integral models
describe holomorphic matrices with a defined ordinary reduction and
infinitesimal corrections. They are not the entire unrestricted matrix
geometry over $\SC$, and our theorems are strong within the former category
without conflating the two.

\section{Positive rigidity on Stein manifolds}\label{nab:sec:stein}

\subsection{An arbitrary-cover theorem with a support hypothesis}

Let $X$ be a connected Stein manifold and $(U_i)_{i\in I}$ a set-indexed open
cover; on disconnected intersections, support statements refer to all connected
components. Suppose
\[
 G_{ij}\in\GG(U_i\cap U_j),\qquad G_{ij}G_{jk}=G_{ik},\quad G_{ii}=I_r.
\]
Call the cocycle \emph{globally support-admissible} if
\begin{equation}\label{nab:eq:Ecover}
 E=\bigcup_{i,j}\supp(G_{ij}-I_r)
\end{equation}
is well ordered. This is a sufficient condition \emph{on a presentation}, not
an invariant necessary condition: regular local gauges may destroy it without
destroying triviality. It must not be conflated with the invariant $\EU(G)$ of
\cref{nab:def:polsupport}, nor with $E_\rho$ of \cref{nab:sec:rh}.

\begin{theorem}[Support-controlled Stein splitting]\label{nab:thm:Stein}
A globally support-admissible positive matrix cocycle on a Stein manifold
splits on its original cover: $G_{ij}=B_iB_j^{-1}$ with $B_i\in\GG(U_i)$. All
supports $\supp(B_i-I_r)$ can be contained in the one monoid
$\angles E\setminus\{0\}$.
\end{theorem}

\begin{proof}
Set $M=\angles E$ and solve $G_{ij}B_j=B_i$ coefficientwise with
$B_i=I_r+\sum_{\gamma\in M\setminus\{0\}}b_{i,\gamma}t^\gamma$. Assume the
equations hold at all exponents below $\gamma$; the coefficient equation at
$\gamma$ is
\begin{equation}\label{nab:eq:Steinrec}
 b_{i,\gamma}-b_{j,\gamma}
 =g_{ij,\gamma}
 +\sum_{\substack{\eta+\theta=\gamma\\ \eta,\theta>0}}g_{ij,\eta}b_{j,\theta}
 =:R_{ij,\gamma},
\end{equation}
a finite sum using only earlier coefficients. The cocycle identity
$G_{ij}G_{jk}=G_{ik}$ together with the equations already solved below
$\gamma$ implies $R_{ij,\gamma}+R_{jk,\gamma}=R_{ik,\gamma}$ on triple
overlaps; expand
$G_{ij}(G_{jk}B_k-B_j)+(G_{ij}B_j-B_i)=G_{ik}B_k-B_i$ at exponent $\gamma$ and
note that all positive multipliers of a lower defect vanish by induction.

Thus $R_{ij,\gamma}$ is an ordinary additive cocycle with values in
$\OO^{r^2}$. Cartan's theorem B gives $H^1(X,\OO^{r^2})=0$
\cite{Cartan,Schmitt}, so it has a coboundary $b_{i,\gamma}$ on the original
cover; this is again the torsor interpretation of first cohomology and assumes
nothing about coherence of the Hahn sheaf. Recursion on $M\setminus\{0\}$
gives the required matrices with a common support bound, and positive
inversion proves they are invertible. No new ordinary domains and no new Hahn
exponents are introduced.
\end{proof}

\subsection{Deformations of a nontrivial ordinary bundle}

The principle is not limited to a trivial ordinary reduction. Let
$\mathcal V$ be an ordinary holomorphic rank-$r$ bundle on $X$ with transition
matrices $g^0_{ij}$, and suppose Hahn transitions $G_{ij}$ have reduction
$g^0_{ij}$ and are invertible over $\II$. The relative transitions
$G_{ij}(g^0_{ij})^{-1}-I_r$ then have positive support on overlaps.

\begin{theorem}[Support-controlled deformation rigidity]
\label{nab:thm:deformation}
If the union of these relative positive supports is well ordered and $X$ is
Stein, there exist positive local matrices $B_i$ with
\begin{equation}\label{nab:eq:deformation}
 G_{ij}=B_ig^0_{ij}B_j^{-1},
\end{equation}
and the nonconstant supports of all $B_i$ lie in the monoid generated by the
given support union. Thus the integral Hahn deformation is isomorphic, by an
isomorphism reducing to the identity, to the scalar extension of its ordinary
bundle.
\end{theorem}

\begin{proof}
Proceed as in \cref{nab:thm:Stein}, now solving $G_{ij}B_j=B_ig^0_{ij}$. At a
fixed positive exponent, multiply the residual equation on the right by
$(g^0_{ij})^{-1}$; the unknown coefficients appear as
$b_{i,\gamma}-g^0_{ij}b_{j,\gamma}(g^0_{ij})^{-1}$. The residual is an
additive cocycle for the ordinary endomorphism sheaf $\Endsh(\mathcal V)$, the
cocycle identity following from $G_{ij}G_{jk}=G_{ik}$ and the lower solved
equations exactly as before. That sheaf is locally free, hence coherent, and
Cartan's theorem B gives $H^1(X,\Endsh(\mathcal V))=0$. Solve the ordinary
coefficient cocycle and continue the support-controlled recursion. All
multiplications by $g^0_{ij}$ and its inverse occur at exponent zero and do not
enlarge the exponent set.
\end{proof}

This is a formal deformation-rigidity argument with a precise support
certificate. It is not a new proof of the ordinary Oka--Grauert principle, and
it is not an unconditional theorem about every Hahn-valued deformation on a
noncompact base. Note that the cohomological input is acyclicity of
$\Endsh(\mathcal V)$, not of $\OO^{r^2}$: the theorem is about deformations of
a genuinely nontrivial ordinary bundle.

\subsection{A cyclic versus noncyclic dichotomy}

\begin{lemma}[Ordered-group alternative]\label{nab:lem:groups}
For a nonzero ordered abelian group $\Gamma$, the positive cone contains no
infinite strictly decreasing sequence if and only if $\Gamma=\Z\delta$ for a
least positive element $\delta$.
\end{lemma}

\begin{proof}
The cyclic case is immediate. Conversely, if there is no least positive
element, choose a smaller positive element repeatedly. If there is a least
positive $\delta$ but $\Gamma\neq\Z\delta$, choose a positive
$\gamma\notin\Z\delta$; then $\gamma-n\delta>0$ for every integer $n\geq0$,
since the first nonpositive difference would either leave an element strictly
between $0$ and $\delta$ or force $\gamma\in\Z\delta$. The sequence
$\gamma,\gamma-\delta,\ldots$ is strictly decreasing and positive.
\end{proof}

\begin{theorem}[Universal positive rigidity and its failure]
\label{nab:thm:dichotomy}
Fix a rank $r\geq1$. Every positive rank-$r$ cocycle on every Stein manifold
is a coboundary if and only if
\[
 \Gamma=\{0\}\quad\text{or}\quad\Gamma=\Z\delta\ \text{for some }\delta>0.
\]
For every noncyclic nonzero $\Gamma$ and every $r\geq2$ there is a nonsplit
determinant-one positive rank-$r$ cocycle already on the puncture-and-disk
cover of $\C$.
\end{theorem}

\begin{proof}
In the cyclic case every transition support is a subset of the universal well
ordered positive set $\{\delta,2\delta,\ldots\}$, so its global union is
automatically well ordered and \cref{nab:thm:Stein} applies on any cover; the
zero group gives only the identity cocycle.

For the converse choose a strictly decreasing positive sequence $\gamma_n$ by
\cref{nab:lem:groups} and set
\[
 g_n=\exp\Bigl(\frac{t^{\gamma_n}}{z-n}\Bigr).
\]
This scalar matrix is its own normalized polar factor, every nonconstant
coefficient being a multiple of a negative power of $z-n$, and its support
union contains the descending set $\{\gamma_n\}$, so it is nonsplit. In rank
$r\geq2$ use
\begin{equation}\label{nab:eq:SLcounter}
 G_n=\operatorname{diag}(g_n,g_n^{-1},1,\ldots,1).
\end{equation}
Every nonconstant coefficient of this matrix is polar, so $P_n=G_n$; the same
descending support prevents splitting in the full positive matrix group, while
$\det G_n=1$. For rank one the scalar example suffices.
\Cref{nab:prop:integraltrivial} converts these examples into nontrivial
integral vector bundles.
\end{proof}

The rank-one part agrees with the Picard-group vanishing dichotomy of the
repository's global-divisor report \cite{RepoGlobal}, which vanishes exactly
for $\Gamma=\{0\}$ or $\Gamma=\Z\delta$; that statement is the rank-one shadow
of \cref{nab:thm:dichotomy}, driven by the same universal set
$\{\delta,2\delta,\ldots\}$, and we cite rather than reprove it. The
higher-rank conclusion here is stated for integral positive deformations, and
the diagonal family \eqref{nab:eq:SLcounter} by itself supplies no nonabelian
information: that is what \cref{nab:thm:hidden} and \cref{nab:thm:depth} are
for.
\section{The two well-ordered-support criteria are not the same criterion}
\label{nab:sec:twocriteria}

Both headline theorems of this report read ``the union of the supports is well
ordered''. The unions concern different objects: full raw transition support
gives only a sufficient gluing test, while raw meridian support gives the exact
monodromy test.

\begin{center}\small
\begin{tabular}{@{}L{0.17\textwidth}L{0.38\textwidth}L{0.38\textwidth}@{}}
\toprule
& \textbf{Criterion P} (\cref{nab:thm:criterion}) &
  \textbf{Criterion M} (\cref{nab:thm:RH})\\
\midrule
Problem &
multiplicative gluing $G_a=B_0B_a^{-1}$ on a puncture-and-disk cover &
realization of a prescribed monodromy by a global positive connection\\[0.3em]
Tested object &
$\Pol(G_a)-I_r$, the normalized polar factor of an \emph{independent local
factorization} of each transition matrix &
$\rho(\ell_d)-I_r$, the prescribed monodromy matrices themselves, for a
compatible based-meridian basis\\[0.3em]
Tested union &
$\EU(G)=\bigcup_a\supp(\Pol(G_a)-I_r)$ &
$E_\rho=\bigcup_{d}\supp(\rho(\ell_d)-I_r)$\\[0.3em]
Raw data &
well-ordering of full raw transition support is sufficient, not necessary;
raw singular support fails both directions
(\cref{nab:thm:hidden,nab:ex:repair}) &
well-ordering of raw meridian support is the invariant condition, independent
of basis and base point by \cref{nab:prop:intrinsic}\\[0.3em]
Must one factor first? &
Yes. Skipping the factorization tests an object that is provably inequivalent
&
No. The monodromy matrices are already the right object; a local
factorization of $\rho(\ell_d)$ is not part of the criterion\\[0.3em]
Invariance &
under arbitrary local holomorphic right gauges
(\cref{nab:cor:localgauges}), coordinate change
(\cref{nab:cor:coordinates}) and ordered value-group extension
(\cref{nab:cor:extension}) &
under change of free meridian basis and base point
(\cref{nab:prop:intrinsic}), one fixed $\GL_r(\K)$ conjugation and ordered
value-group enlargement (\cref{nab:prop:conjugation})\\
\bottomrule
\end{tabular}
\end{center}

\medskip
The trap is concrete, and it has already been printed. In
\cref{nab:thm:hidden} the raw singular support is the single exponent $\{1\}$
and no splitting exists; in \cref{nab:ex:repair}, \emph{adding} a term whose
exponents form the descending set $\{1+1/n\}$ among the raw poles
\emph{removes} the obstruction. So for Criterion P the raw singular support is
not the invariant, and the factorization must be performed first. Nothing of
the kind happens for Criterion M: there the union is over the prescribed
matrices, \cref{nab:prop:intrinsic} shows the condition does not depend on the
presentation at all, and \cref{nab:ex:descending} obstructs precisely because
the prescribed matrices themselves carry $\{1/n\}$.

Two further consequences of the difference are worth stating so that neither
criterion is quoted for the other's conclusion. First, the objects tested live
in different places: $\Pol(G_a)$ is a matrix of negative Laurent tails in a
local coordinate at $a$, whereas $\rho(\ell_d)$ is a matrix of Hahn constants.
Second, the two criteria are not two instances of one general theorem proved
here; they are two theorems with two proofs, sharing
\cref{nab:lem:neumann,nab:lem:inverse} and an ordinary Mittag--Leffler input.
This report does not claim a common generalization, and
\cref{nab:sec:open} lists the obvious attempt as open.
\part*{Part II. Inverse monodromy: the meridian-support criterion}
\addcontentsline{toc}{section}{\textbf{Part II. Inverse monodromy: the meridian-support criterion}}

\section{Connections, fundamental matrices, and holonomy}
\label{nab:sec:holonomy}

A positive connection on the trivial rank-$r$ bundle over an ordinary domain
$U$ means
\begin{equation}\label{nab:eq:connection}
 dY=\Omega Y,\qquad
 \Omega=\sum_{\gamma\in S}t^\gamma\Omega_\gamma,
 \qquad S\subseteq\Gamma_{>0}\ \text{well ordered},
\end{equation}
where each $\Omega_\gamma$ is an ordinary holomorphic matrix one-form on
\emph{all} of $U$. In complex dimension one such connections are automatically
flat: both the exterior derivative of a holomorphic one-form and the wedge of
two of them vanish. As always, ``positive'' concerns exponents only.

Let $\widetilde U\to U$ be the ordinary universal covering with a chosen lift
$\widetilde p$ of the base point $p$; pullbacks of coefficients are
understood.

\begin{theorem}[Fundamental matrix with a support certificate]
\label{nab:thm:fundamental}
For a connection \eqref{nab:eq:connection} there is a unique positive
fundamental matrix
\[
 Y=I_r+\sum_{\gamma>0}t^\gamma Y_\gamma\ \text{on }\widetilde U,
 \qquad dY=\Omega Y,\quad Y(\widetilde p)=I_r.
\]
Every coefficient is holomorphic on all of $\widetilde U$ and
\[
 \supp(Y-I_r)\cup\supp(Y^{-1}-I_r)\subseteq\angles S\setminus\{0\}.
\]
Analytic continuation determines a representation
$\Hol_\Omega:\pi_1(U,p)\to I_r+\Mat_r(\mm)$ with the same uniform support
bound for \emph{every} group element.
\end{theorem}

\begin{proof}
Write $T=\angles S\setminus\{0\}$ and define the coefficients by well-founded
recursion on $T$, missing coefficients of $\Omega$ being zero. At weight
$\gamma$ set
\begin{equation}\label{nab:eq:solrec}
 dY_\gamma=\Omega_\gamma
 +\sum_{\substack{\alpha+\beta=\gamma\\ \alpha,\beta>0}}\Omega_\alpha Y_\beta,
 \qquad Y_\gamma(\widetilde p)=0.
\end{equation}
Only finitely many pairs contribute by \cref{nab:lem:neumann}, and
$\beta<\gamma$. The right-hand side is a holomorphic matrix one-form on the
simply connected Riemann surface $\widetilde U$, so its entries have unique
holomorphic primitives vanishing at $\widetilde p$; this constructs $Y_\gamma$
on the whole cover. The family has support in the well-ordered set $T$ and
satisfies the equation coefficientwise, and \cref{nab:lem:inverse} gives its
inverse with the same monoid support.

For uniqueness even against a solution with a different positive well-ordered
support, take the least exponent at which two solutions differ, in the
well-ordered union of their supports: at that exponent the difference has zero
derivative and zero value at $\widetilde p$, hence is zero, a contradiction.

A continued fundamental matrix differs from the original branch on a connected
overlap by a constant matrix on the right, since $d(Y^{-1}\widehat Y)=0$
follows coefficientwise from the equation. Continuing around a based loop gives
the normalized monodromy matrix, whose coefficients are values of coefficients
of $Y$; therefore all its positive exponents lie in $T$, independently of the
loop. Uniqueness of continuation gives the group law with the convention of
\cref{nab:lem:free}, and homotopy invariance follows from the universal-cover
construction.
\end{proof}

This theorem is the precise form of the \emph{forward} map that the
repository's differential-equations report constructs from a coherent
connection, recording that the exponent-zero coefficient is the ordinary
monodromy matrix and that all other exponents lie in the generated monoid
\cite{RepoDiff}. Here the exponent-zero coefficient is $I_r$ by positivity.
\Cref{nab:thm:RH} is the converse.

\subsection{The explicit iterated-integral expression}

For a piecewise smooth path $c:[0,1]\to U$ write
$c^*\Omega_\gamma=A_\gamma(s)\,ds$. Its transport is
\begin{equation}\label{nab:eq:dyson}
 T_\Omega(c)=I_r+\sum_{n\geq1}\ \sum_{\gamma_1,\ldots,\gamma_n\in S}
 t^{\gamma_1+\cdots+\gamma_n}
 \int_{0<s_1<\cdots<s_n<1}
 A_{\gamma_n}(s_n)\cdots A_{\gamma_1}(s_1)\,ds_1\cdots ds_n.
\end{equation}
This is the usual ordered-integral solution, a classical mechanism related to
iterated path integrals \cite{Chen}; the name ``Dyson expansion'' is used for
it purely as classical mathematics, with no physical reading attached. Its
meaning here is strict: the outer family is strongly summable by the
finite-word clause of \cref{nab:lem:neumann}, and each individual coefficient
is a \emph{finite} sum of ordinary iterated integrals. The expression agrees
with \cref{nab:thm:fundamental}, since termwise differentiation gives
\eqref{nab:eq:solrec}; for a loop it gives $\Hol_\Omega$.

\begin{lemma}[Triangular change of monodromy]\label{nab:lem:triangular}
Suppose two positive connections $\Omega$ and $\Psi$ agree at every exponent
strictly below $\gamma$. Then their transports along any fixed path agree below
$\gamma$, and for a based loop $\ell$,
\begin{equation}\label{nab:eq:triangular}
 \coeff{\gamma}\bigl(\Hol_\Psi(\ell)-\Hol_\Omega(\ell)\bigr)
 =\int_\ell(\Psi_\gamma-\Omega_\gamma).
\end{equation}
Terms strictly above $\gamma$ cannot affect coefficients at or below $\gamma$.
\end{lemma}

\begin{proof}
Compare the word expansions \eqref{nab:eq:dyson} over the well-ordered union
of the two positive supports. A differing coefficient of weight at least
$\gamma$ cannot occur in a word of total weight less than $\gamma$; at weight
exactly $\gamma$ it can occur only as the single letter of the word, every
additional letter having strictly positive weight. The length-one term is the
displayed ordinary period. All comparisons are finite at each exponent.
\end{proof}

The crucial feature is that the linearized map at \emph{every} successive
weight is the same ordinary period map $P$ of \eqref{nab:eq:periodmap}.
Noncommutativity appears in the already determined lower-weight terms; it does
not change the operator needed to solve the new correction equation.

\section{The exact realization theorem}\label{nab:sec:rh}

\begin{definition}[Uniform support admissibility and $\sigma$-normalization]
\label{nab:def:admissible}
For $\rho:\pi_1(U,p)\to I_r+\Mat_r(\mm)$ with $U=\C\setminus D$ and
$M_d=\rho(\ell_d)$, put
\begin{equation}\label{nab:eq:Erho}
 E_\rho=\bigcup_{d\in D}\supp(M_d-I_r),
\end{equation}
the union being over all matrix entries and all based meridians of the fixed
free basis of \cref{nab:lem:free}. The representation is \emph{uniformly
support-admissible} when $E_\rho$ is well ordered. For a fixed period section
$\sigma$ of \cref{nab:lem:period}, a connection is
\emph{$\sigma$-normalized} if every entry of every coefficient lies in
$\im\sigma\subseteq\LL(D)$.
\end{definition}

\begin{theorem}[Support-sensitive Riemann--Hilbert realization: Criterion M]
\label{nab:thm:RH}
Let $D\subset\C$ be closed and discrete, $p\in U=\C\setminus D$, with free
based meridians fixed. For $\rho:\pi_1(U,p)\to I_r+\Mat_r(\mm)$ the following
are equivalent:
\begin{enumerate}[label=\textup{(\roman*)}]
\item $E_\rho$ is well ordered;
\item $\rho$ is the monodromy of a positive globally common-domain Hahn
connection on $U$;
\item $\rho$ is the monodromy of such a connection whose coefficients are
meromorphic on $\C$ with at most simple poles in $D$ and no other finite
poles.
\end{enumerate}
If these hold then for every fixed complex-linear period section $\sigma$
there is exactly one $\sigma$-normalized realizer $\Omega^\sigma$, and
\begin{equation}\label{nab:eq:RHsupport}
 \supp\Omega^\sigma\subseteq\angles{E_\rho}\setminus\{0\}.
\end{equation}
Uniqueness is among \emph{all} positive globally supported
$\sigma$-normalized connections, not only those whose support was assumed in
advance to lie in the right-hand side of \eqref{nab:eq:RHsupport}.
\end{theorem}

\begin{proof}
\emph{Necessity.} If $\rho=\Hol_\Omega$ with $\supp\Omega=S$ then
\cref{nab:thm:fundamental} gives
$\supp(\rho(g)-I_r)\subseteq\angles S\setminus\{0\}$ for \emph{every} group
element $g$, so $E_\rho$ is a subset of one well-ordered set. This proves
(ii)$\Rightarrow$(i), and (iii)$\Rightarrow$(ii) is immediate.

\emph{Construction.} Assume (i), put $T=\angles{E_\rho}\setminus\{0\}$, and
write $M_d=I_r+\sum_{\gamma\in T}t^\gamma m_{d,\gamma}$, missing coefficients
being zero. We construct $\Omega_\gamma\in\Mat_r(\LL(D))$ for $\gamma\in T$ by
recursion. Suppose all lower coefficients have been constructed and form the
legitimate positive connection
$\Omega^{<\gamma}=\sum_{\delta\in T,\ \delta<\gamma}t^\delta\Omega_\delta$,
whose support is a well-ordered subset of $T$; the initial segment need not be
finite. Define the lower-order prediction and the remaining defect by
\begin{equation}\label{nab:eq:defect}
 h^{<\gamma}_{d,\gamma}=\coeff{\gamma}\Hol_{\Omega^{<\gamma}}(\ell_d),
 \qquad
 b_{d,\gamma}=m_{d,\gamma}-h^{<\gamma}_{d,\gamma}.
\end{equation}
The prediction exists by \cref{nab:thm:fundamental}; equivalently it is the
finite sum of all iterated-integral words of total weight $\gamma$ formed from
previously constructed coefficients. Apply the period section entrywise and
set
\begin{equation}\label{nab:eq:RHrecursion}
 \Omega_\gamma=\sigma\bigl((b_{d,\gamma})_{d\in D}\bigr).
\end{equation}
The datum at one exponent is an arbitrary matrix-valued element of the
\emph{product} $\C^D$, and \cref{nab:lem:period} applies even when infinitely
many of its entries are nonzero; this is why no summability across punctures is
required.

The coefficient $\Omega_\gamma$ is meromorphic on $\C$ with at most simple
poles in $D$. It corrects exactly the defect at weight $\gamma$, since
$\int_{\ell_d}\Omega_\gamma=b_{d,\gamma}$ and \cref{nab:lem:triangular}
applies, and it leaves all lower weights unchanged. Transfinite recursion on
the well-ordered set $T$ therefore constructs
$\Omega^\sigma=\sum_{\gamma\in T}t^\gamma\Omega_\gamma$, all of whose
coefficients are holomorphic on the same domain $U$ and whose support is a
subset of $T$. No limiting assertion at a limit ordinal is involved: a
coefficient is simply assigned to every member of a well-ordered set.

To verify the monodromy, fix $\gamma\in T$. Coefficients at weights above
$\gamma$ cannot affect it by \cref{nab:lem:triangular}, and the recursion has
already set it equal to $m_{d,\gamma}$. The monodromy support of the final
connection lies in $\angles T\setminus\{0\}=T$, so no exponent outside $T$
appears. Thus every meridian has its prescribed monodromy, and the meridians
freely generate the fundamental group by \cref{nab:lem:free}, proving
$\Hol_{\Omega^\sigma}=\rho$ and (iii).

\emph{Uniqueness.} Let $\Omega$ and $\Psi$ be two positive
$\sigma$-normalized realizers. If they differ, the union of their supports is
well ordered and there is a least differing exponent $\gamma$. Their
holonomies are equal, so \cref{nab:lem:triangular} gives
$P(\Psi_\gamma-\Omega_\gamma)=0$. But $\Psi_\gamma-\Omega_\gamma$ lies
entrywise in $\im\sigma$, and $P\sigma=\id$ says that $P$ is injective there.
Hence the difference is zero, a contradiction.
\end{proof}

\begin{remark}[No hidden convergence over the puncture set]
\label{nab:rem:noconv}
The construction is not a sum of one independently built differential equation
per puncture. At a fixed Hahn exponent it first collapses \emph{all} puncture
data into one ordinary meromorphic one-form using the period section, and only
then are these ordinary forms assembled in Hahn exponents. This move is
indispensable when infinitely many punctures use the same exponent; see
\cref{nab:ex:sameexponent}.
\end{remark}

\subsection{An intrinsic version of the support condition}

\begin{proposition}[Independence of the meridian basis]
\label{nab:prop:intrinsic}
For a near-identity representation $\rho$ put
$E_{\mathrm{all}}(\rho)=\bigcup_{g\in\pi_1(U,p)}\supp(\rho(g)-I_r)$. Then
$E_\rho$ is well ordered if and only if $E_{\mathrm{all}}(\rho)$ is well
ordered. Hence admissibility does not depend on the chosen free meridians or,
after a base-point identification, on the base point.
\end{proposition}

\begin{proof}
One implication follows from $E_\rho\subseteq E_{\mathrm{all}}(\rho)$. For the
other, assume $E_\rho$ well ordered. Each group element is a finite word in the
$M_d$ and their inverses; \cref{nab:lem:inverse} expands the inverses in the
common monoid $\angles{E_\rho}$, and multiplication preserves that monoid, so
$E_{\mathrm{all}}(\rho)\subseteq\angles{E_\rho}\setminus\{0\}$. Changing based
meridians, or identifying fundamental groups at different base points, changes
the presentation of the same representation and not this condition on all its
values.
\end{proof}

The support set $E_\rho$ itself can change with the meridian basis; its
\emph{well-ordering} is the invariant proved here. Contrast
\cref{nab:sec:twocriteria}: Criterion P is \emph{not} a
condition on the raw transition data, and the analogue of
\cref{nab:prop:intrinsic} for $\bigcup_a\supp(G_a-I_r)$ is false.

\begin{proposition}[Fixed conjugation and enlargement cannot repair support]
\label{nab:prop:conjugation}
For a family $(M_d)$ over $\K$, well-ordering of
$\bigcup_d\supp(M_d-I_r)$ is invariant under conjugation by any one fixed
$C\in\GL_r(\K)$, and is unchanged when the family is embedded by an
order-preserving value-group embedding into a larger Hahn workspace. An
arbitrary conjugation need not preserve near-identity positivity; the assertion
about well-ordered supports does not require that property.
\end{proposition}

\begin{proof}
Let $A$ and $B$ be the supports of the finitely many entries of $C$ and
$C^{-1}$; they are well ordered. If $E$ is the union of the supports of
$M_d-I_r$ then
$\bigcup_d\supp(C(M_d-I_r)C^{-1})\subseteq A+E+B$, which is well ordered by
the last clause of \cref{nab:lem:neumann} when $E$ is. The converse follows by
conjugating with $C^{-1}$. An order embedding preserves and reflects the
ordering of a subset, hence preserves and reflects its being well ordered.
\end{proof}

\Cref{nab:prop:conjugation} for Criterion M and \cref{nab:cor:extension} for
Criterion P are two independent proofs, for two different objects, of the same
moral: enlarging the surreal workspace never repairs these obstructions.

\section{Framed gauge classification and removable gauges}
\label{nab:sec:gauge}

\begin{definition}[Positive framed gauge]\label{nab:def:gauge}
A positive gauge on $U$ is a matrix $G\in\GG(U)$. It is \emph{framed} at $p$
when $G(p)=I_r$ as an equality of Hahn matrices. Its inverse is a positive
globally supported holomorphic matrix by \cref{nab:lem:inverse}, and it
transforms a connection by
\begin{equation}\label{nab:eq:gaugeact}
 \Psi=dG\,G^{-1}+G\Omega G^{-1}.
\end{equation}
\end{definition}

\begin{theorem}[Framed classification with controlled support]
\label{nab:thm:gauge}
Two positive globally Hahn-supported connections $\Omega$ and $\Psi$ on $U$
have the same based monodromy if and only if there is a positive framed gauge
$G$ satisfying \eqref{nab:eq:gaugeact}. Such a gauge is unique, and if
$E=\supp\Omega\cup\supp\Psi$ then
\begin{equation}\label{nab:eq:gaugesupport}
 \supp(G-I_r)\cup\supp(G^{-1}-I_r)\subseteq\angles E\setminus\{0\}.
\end{equation}
Consequently holonomy is a bijection between framed positive gauge classes and
uniformly support-admissible near-identity representations.
\end{theorem}

\begin{proof}
Let $Y_\Omega,Y_\Psi$ be the normalized fundamental matrices on
$\widetilde U$. If the based monodromies agree, form
$\widetilde G=Y_\Psi Y_\Omega^{-1}$. Upon continuation around any loop both
fundamental matrices acquire the same right factor $M$, so the product becomes
$Y_\Psi M(Y_\Omega M)^{-1}=Y_\Psi Y_\Omega^{-1}$. Each ordinary coefficient
descends to $U$, and the whole family already has support in $\angles E$ by
\cref{nab:thm:fundamental,nab:lem:inverse}. Thus the descended matrix is a
globally supported positive gauge $G$ with $G(p)=I_r$, and differentiation
gives $dG=\Psi G-G\Omega$, equivalent to \eqref{nab:eq:gaugeact}.

Conversely, if a framed gauge exists then $GY_\Omega$ solves the $\Psi$
equation with initial value $I_r$, so uniqueness in
\cref{nab:thm:fundamental} gives $GY_\Omega=Y_\Psi$. Since $G$ is single
valued on $U$ the monodromy matrices agree, and the equality also forces
$G=Y_\Psi Y_\Omega^{-1}$, proving uniqueness and the support bound. The
bijection now follows from \cref{nab:thm:RH}.
\end{proof}

The logarithmic subclass has a stronger rigidity property: although the
comparison gauge is initially defined only on the punctured plane, no
coefficient of it actually needs a puncture.

\begin{theorem}[Entire extension of a logarithmic comparison gauge]
\label{nab:thm:entiregauge}
Suppose $\Omega$ and $\Psi$ are as in \cref{nab:thm:gauge}, have the same
based monodromy, and every coefficient of both lies in $\Mat_r(\LL(D))$. Then
the unique framed gauge and its inverse extend coefficientwise holomorphically
to all of $\C$; in particular $G,G^{-1}\in\GG(\C)$, with the same support
bound \eqref{nab:eq:gaugesupport}.
\end{theorem}

\begin{proof}
Write $G=I_r+\sum_{\gamma>0}t^\gamma G_\gamma$. The gauge equation at weight
$\gamma$ reads
\begin{equation}\label{nab:eq:gaugerec}
 dG_\gamma=\Psi_\gamma-\Omega_\gamma
 +\sum_{\substack{\alpha+\beta=\gamma\\ \alpha,\beta>0}}
  (\Psi_\alpha G_\beta-G_\beta\Omega_\alpha).
\end{equation}
Induct on the positive support monoid in \eqref{nab:eq:gaugesupport}. If every
lower $G_\beta$ is entire, the right-hand side is an ordinary meromorphic
one-form on $\C$ with at most simple poles in $D$, the sum being finite. Fix
$d\in D$ and a disk containing no other point of $D$. The coefficient
$G_\gamma$ is already single valued and holomorphic on the punctured disk, so
the integral of its derivative around a small circle vanishes; hence the
right-hand side of \eqref{nab:eq:gaugerec} has residue zero at $d$, and having
at most a simple pole it is holomorphic there. Its primitive on the full disk
differs from $G_\gamma$ by a constant on the punctured disk, so $G_\gamma$
extends holomorphically across $d$. Doing this for all $d$ extends $G_\gamma$
to an entire function, with no uniform bound on the sizes of coefficient
functions required, and the ordinary domain after extension is the same $\C$
for every coefficient. Finally the geometric inverse of a positive entire
family has entire coefficients and the same monoid support.
\end{proof}

\begin{corollary}[A global normal slice for logarithmic connections]
\label{nab:cor:slice}
For a fixed $\sigma$, every positive logarithmic connection has exactly one
$\sigma$-normalized representative under framed positive entire gauges, namely
the connection constructed from its monodromy by
\eqref{nab:eq:RHrecursion}.
\end{corollary}

\begin{proof}
Its monodromy is admissible by \cref{nab:thm:fundamental};
\cref{nab:thm:RH,nab:thm:entiregauge} give existence, and uniqueness of the
normalized connection and of the framed gauge has already been proved.
\end{proof}

\begin{example}[A nonlogarithmic connection with trivial monodromy]
\label{nab:ex:nonlog}
On $\C^*$ consider the scalar connection $\Omega=-t\,dz/z^2$ with base point
$p=1$. Its normalized solution is $Y=\exp(t(1/z-1))$, a positive globally
supported holomorphic family on $\C^*$, single valued, so the monodromy is
trivial. The unique framed gauge from $\Omega$ to the zero connection is
$G=Y^{-1}=\exp(-t(1/z-1))$, whose coefficient at $t$ has a pole at $0$.
Therefore \cref{nab:thm:entiregauge} cannot be extended to connections with
arbitrary higher-order coefficient poles.
\end{example}

\section{Finite-puncture inversion and determinant one}\label{nab:sec:sl}

\begin{corollary}[Unique positive residue matrices for finite data]
\label{nab:cor:finite}
Let $D=\{a_1,\ldots,a_m\}$ be finite, fix based meridians, and let
$M_j\in I_r+\Mat_r(\mm)$ be arbitrary. There are unique residue matrices
$R_j\in\Mat_r(\mm)$ for which
\begin{equation}\label{nab:eq:fuchsian}
 \Omega=\sum_{j=1}^mR_j\frac{dz}{z-a_j}
\end{equation}
has based meridian monodromies $M_j$. If $E=\bigcup_j\supp(M_j-I_r)$ then
$\supp R_j\subseteq\angles E\setminus\{0\}$. The form has at most a simple pole
at infinity, with residue $-\sum_jR_j$ there.
\end{corollary}

\begin{proof}
A finite union of well-ordered sets in a linear order is well ordered, so $E$
is admissible. Use the explicit section \eqref{nab:eq:finitesection} in
\cref{nab:thm:RH}: being normalized for that section is precisely the form
\eqref{nab:eq:fuchsian}, and coefficientwise uniqueness gives uniqueness of all
the $R_j$. In the coordinate $w=1/z$ one has
$dz/(z-a_j)=-dw/(w(1-a_jw))$, proving the assertion at infinity.
\end{proof}

This is a near-identity Hahn version of a classical Fuchsian inverse problem,
not a realization result for arbitrary complex monodromy on a globally trivial
bundle. The residue matrices must have positive support, a regular singularity
at infinity is permitted, no assertion is made that the connection is
holomorphic at infinity, and uniqueness is not claimed among arbitrary residue
matrices with the same monodromy: ordinary integer shifts of exponents, which
already destroy such unrestricted uniqueness, are excluded by the
positive-support hypothesis.

\begin{lemma}[Determinants and positive connections]\label{nab:lem:det}
For the normalized fundamental matrix $Y$ of a positive connection,
$d(\det Y)=(\tr\Omega)\det Y$. Consequently a traceless positive connection
has determinant-one transport along every path.
\end{lemma}

\begin{proof}
The determinant is a finite polynomial in the entries of $Y$, so the adjugate
identity and the Leibniz rule apply coefficientwise, and $Y$ is invertible by
\cref{nab:thm:fundamental}. Jacobi's formula reduces to the displayed
identity. If $\tr\Omega=0$ the determinant is constant along the normalized
solution and equals $1$ at the starting point. All identities involve finite
coefficientwise algebra or already justified support products.
\end{proof}

\begin{theorem}[Traceless normalized realization]\label{nab:thm:SL}
Suppose an admissible representation in \cref{nab:thm:RH} takes values in
$\SL_r(\K)$. For an entrywise complex-linear period section its unique
normalized realizer satisfies $\tr\Omega^\sigma=0$. Conversely every traceless
positive connection has monodromy in $\SL_r(\K)$, and the unique framed gauge
between two traceless connections with the same monodromy has determinant one.
These statements retain all the support and logarithmic-extension conclusions
above.
\end{theorem}

\begin{proof}
Use the recursion \eqref{nab:eq:RHrecursion} and assume every lower
coefficient $\Omega_\delta$ traceless. Then $\Omega^{<\gamma}$ is traceless, so
its monodromy $N_d$ has determinant one by \cref{nab:lem:det}; by construction
$N_d$ and the target $M_d$ agree below $\gamma$, and both have constant term
$I_r$. In the determinant polynomial a difference at weight $\gamma$ can
contribute to that weight only with all other factors of weight zero, so
\[
 \coeff{\gamma}(\det M_d-\det N_d)
 =\tr\bigl(m_{d,\gamma}-h^{<\gamma}_{d,\gamma}\bigr).
\]
The left side is zero, hence the defect vector in \eqref{nab:eq:defect} is
traceless at every puncture. Since the same linear section is used for each
matrix entry,
$\tr\Omega_\gamma=\sigma((\tr b_{d,\gamma})_{d\in D})=0$, proving the
induction. The converse is \cref{nab:lem:det}. For a gauge satisfying
$dG=\Psi G-G\Omega$, Jacobi's formula gives
$d\det G=(\tr\Psi-\tr\Omega)\det G=0$, and its value at $p$ is $1$.
\end{proof}

\section{Sharp global obstructions and calibrating examples}
\label{nab:sec:examples}

\subsection{Finite solvability does not imply global solvability}

\begin{example}[A countable special-linear obstruction]
\label{nab:ex:descending}
Let $D=\{1,2,3,\ldots\}$, $\Gamma=\Q$, and
\[
 N=\begin{pmatrix}0&1\\0&0\end{pmatrix},\qquad M_n=I_2+t^{1/n}N.
\]
Because the meridians form a free basis, this prescription defines a
representation $\rho:\pi_1(\C\setminus D,p)\to\SL_2(\K)$. Each generator matrix
has one positive exponent, is invertible, and is an entirely legitimate Hahn
matrix, and every finite initial collection of the prescribed monodromies is
realizable on the corresponding finitely punctured plane by
\cref{nab:cor:finite}. Nevertheless no positive globally Hahn-supported
holomorphic connection on $\C\setminus D$ has all the prescribed monodromies.
\end{example}

\begin{proof}
The determinant is $1$ because $N^2=0$ and $N$ is strictly upper triangular. A
finite set of the exponents $1/n$ is well ordered, but the full union
$E_\rho=\{1,1/2,1/3,\ldots\}$ is a strictly decreasing sequence with no least
member, hence not well ordered. Nonexistence follows from
\cref{nab:thm:RH}, even without any requirement of simple poles, and
\cref{nab:prop:conjugation} shows that a fixed change of matrix basis or an
enlargement of the value group cannot repair it.
\end{proof}

One-puncture realizers can even be written down:
\[
 \Omega_n=\frac{t^{1/n}}{2\pi i}N\frac{dz}{z-n},
 \qquad
 \exp\Bigl(\int_{\ell_n}\Omega_n\Bigr)=I_2+t^{1/n}N.
\]
The failure is not an inability to solve a differential equation at any one
puncture; it is the impossibility of placing the entire prescription in one
global support set. This is a genuine local-to-global obstruction in the chosen
coefficient category, not a claim that no ordinary local system exists: the
representation itself already defines an algebraic local system over $\K$.

\subsection{The union condition must not be strengthened to point-finiteness}

\begin{example}[Infinitely many punctures at one exponent]
\label{nab:ex:sameexponent}
Let $D=\Z$, $U=\C\setminus\Z$, $\Gamma=\Q$, and use the same nilpotent $N$.
Then
\begin{equation}\label{nab:eq:cot}
 \Omega=\frac{t}{2i}N\cot(\pi z)\,dz
\end{equation}
has monodromy $I_2+tN$ at every positive based meridian. The family
$(M_n-I_2)_{n\in\Z}$ is \emph{not} strongly summable as a family indexed by
punctures, because exponent $1$ appears infinitely many times; it is
nevertheless globally realizable.
\end{example}

\begin{proof}
The ordinary form $\cot(\pi z)\,dz$ has residue $1/\pi$ at every integer, so
each meridian period of \eqref{nab:eq:cot} is $tN$. Every product of two
coefficient matrices vanishes, both being multiples of $N$ with $N^2=0$, so the
Dyson expansion \eqref{nab:eq:dyson} terminates after its length-one term and
gives $I_2+tN$. The connection has the singleton support $\{1\}$ and is
holomorphic on one common domain $U$. Nothing in the definition asks that its
periods over \emph{different} loops form a summable family.
\end{proof}

Together, \cref{nab:ex:descending,nab:ex:sameexponent} identify the precise
strength of Criterion M. Requiring only that individual matrices be legitimate
Hahn matrices is too weak; requiring a strongly summable family over punctures
is too strong; a well-ordered \emph{union} of supports is exactly right.

\subsection{Admissible supports need not be locally finite in an interval}

\begin{example}[Infinitely many exponents below one bound]
\label{nab:ex:boundedsupport}
On $D=\{1,2,\ldots\}$ prescribe $M_n=I_2+t^{2-1/n}N$. The support
$E=\{2-1/n:n\geq1\}$ is increasing of order type $\omega$, hence well ordered,
although it has infinitely many elements below $2$. These data are realizable:
for a fixed period section put $\alpha_n=\sigma(e_n)$ and
\[
 \Omega=N\sum_{n\geq1}t^{2-1/n}\alpha_n,
\]
where $e_n\in\C^D$ is $1$ at $n$ and $0$ elsewhere. The connection has support
in $E$, and its holonomy is exactly the prescribed family because $N^2=0$.
\end{example}

The displayed sum is a Hahn family, not an assertion that its ordinary partial
sums converge in the valuation topology. In $\Q$ the tail valuations $2-1/n$
stay below $2$, so those partial sums do not even approach the full family in
that topology. The example also shows why our theorems supply no procedure
enumerating all coefficients below a fixed bound in finitely many steps.

\section{A noncommuting calculation across two valuation ranks}
\label{nab:sec:mixed}

\subsection{The targets and the first-order guess}

Take $U=\C\setminus\{0,1\}$ and base point $b=1/2$, with the usual positive
based meridians: a real tail to a small circle about the puncture, a positive
circuit, and the reversed tail. Define
\[
 N_+=\begin{pmatrix}0&1\\0&0\end{pmatrix},\quad
 N_-=\begin{pmatrix}0&0\\1&0\end{pmatrix},\quad
 H=[N_+,N_-]=\begin{pmatrix}1&0\\0&-1\end{pmatrix},
\]
so $N_+^2=N_-^2=0$. Initially use two independent formal bookkeeping variables
$u,v$ of positive degree and prescribe
\begin{equation}\label{nab:eq:twotargets}
 M_0=I_2+uN_+,\qquad M_1=I_2+vN_-.
\end{equation}
At the end set $u=t^\alpha$ and $v=t^\beta$ for positive, rationally
independent $\alpha,\beta\in\Gamma$; the choice $\alpha=1$, $\beta=\omega$ is
the rank-two surcomplex example. Put
\[
 \eta_0=\frac{dz}{2\pi iz},\qquad
 \eta_1=\frac{dz}{2\pi i(z-1)},\qquad
 C=\frac{\log2}{2\pi i}.
\]
The obvious first-order guess is
\begin{equation}\label{nab:eq:linearguess}
 \Omega^{(1)}=uN_+\eta_0+vN_-\eta_1.
\end{equation}
Its first-order periods are correct. Its mixed-order monodromy is not.

\begin{proposition}[The first nonabelian correction]\label{nab:prop:mixed}
For the connection \eqref{nab:eq:linearguess} the coefficient of $uv$ in the
based monodromy is $CH$ at the meridian of $0$ and $-CH$ at the meridian of
$1$. Consequently the unique normalized realizer has expansion
\begin{equation}\label{nab:eq:mixedconnection}
 \Omega=uN_+\eta_0+vN_-\eta_1-uv\,CH(\eta_0-\eta_1)+O_{\deg\geq3},
\end{equation}
and its coefficients of $u^2$ and $v^2$ are zero. The remainder notation refers
only to total degree in $u,v$, not to a claimed valuation bound.
\end{proposition}

\begin{proof}
Along a chosen meridian let $L_j(s)=\int_0^s\eta_j$ be the accumulated
integral from the base point. The first-order transport is
$I_2+uN_+L_0+vN_-L_1$, so the mixed term in the coefficient equation is
\begin{equation}\label{nab:eq:mixedintegral}
 N_+N_-\int_\ell L_1\eta_0+N_-N_+\int_\ell L_0\eta_1.
\end{equation}
For the meridian of $0$, $L_1$ is the single-valued ordinary function
$L_1(z)=(\log(1-z)-\log(1-b))/2\pi i$ on a neighbourhood containing the tail
and the disk about $0$ but not $1$; it satisfies $L_1(b)=0$ and $L_1(0)=C$.
Integrating $d(L_0L_1)$ along the full loop gives
$\int_\ell L_0\eta_1=-\int_\ell L_1\eta_0$, because the product vanishes at
both endpoints, $L_0$ ending at $1$ but $L_1$ ending at $0$. The residue
theorem gives $\int_\ell L_1\eta_0=L_1(0)=C$, contributions of the outgoing and
incoming tails cancelling. Thus \eqref{nab:eq:mixedintegral} equals
$C(N_+N_--N_-N_+)=CH$.

For the meridian of $1$ use instead $L_0(z)=(\log z-\log b)/2\pi i$, which is
single valued near that tail and disk with $L_0(b)=0$ and $L_0(1)=C$; now
$\int_\ell L_0\eta_1=C$ and $\int_\ell L_1\eta_0=-C$, so the mixed coefficient
is $-CH$. The prescribed matrices \eqref{nab:eq:twotargets} have zero $uv$
coefficient, so the ordinary period section for $\{0,1\}$ assigns the
correcting one-form $-CH\eta_0+CH\eta_1$ at that degree, which is
\eqref{nab:eq:mixedconnection}. Pure second-order terms in the uncorrected
transport vanish because $N_+^2=N_-^2=0$, so their defects and normalized
correction coefficients are zero.
\end{proof}

In residue form the same calculation reads
\begin{align}
 R_0&=\frac{u}{2\pi i}N_+-\frac{uv\log2}{(2\pi i)^2}H+O_{\deg\geq3},
 \label{nab:eq:R0}\\
 R_1&=\frac{v}{2\pi i}N_-+\frac{uv\log2}{(2\pi i)^2}H+O_{\deg\geq3}.
 \label{nab:eq:R1}
\end{align}
The mixed corrections cancel in the residue at infinity but are essential at
the two finite poles. All displayed coefficient matrices are traceless, in
agreement with \cref{nab:thm:SL}.

\subsection{The representation is genuinely nonabelian}

The exact group commutator of the prescribed matrices is
\begin{equation}\label{nab:eq:commutator}
 M_0M_1M_0^{-1}M_1^{-1}
 =\begin{pmatrix}1+uv+u^2v^2&-u^2v\\ uv^2&1-uv\end{pmatrix},
\end{equation}
with determinant $1$ and leading mixed term $uvH\neq0$. So the realization
theorem here is not an entrywise construction of independent scalar equations:
the extra commutator in \eqref{nab:eq:mixedconnection} is exactly the
interaction that a scalar period prescription would miss. This is the Part~II
counterpart of the cross term $-a_-c_+$ of \cref{nab:prop:heisenberg}: both
constructions require mixed products absent from the abelianized equations.
Here the mixed term is a solvable correction; in Part~I such products can also
create a globally inadmissible polar-support union.

\subsection{Specialization to a genuinely higher-rank Hahn group}

Set $\Gamma=\Q+\Q\omega$, $u=t$ and $v=t^\omega$. The monoid generated by $1$
and $\omega$ is $\{a+b\omega:a,b\in\N\cup\{0\}\}$, ordered first by $b$ and
then by $a$, and it is well ordered. The correction in
\eqref{nab:eq:mixedconnection} occurs at exponent $1+\omega$. No finite number
of increments by $1$ reaches even $\omega$, but the well-founded recursion
handles $\omega$ and $1+\omega$ directly, with no change in the proof or in its
support certificate; compare \cref{nab:rem:nonarch}.

The notation $O_{\deg\geq3}$ is especially important here. A term $t^3$ has
total degree three and yet smaller valuation than $t^\omega$, while
$t^{3\omega}$ has much larger valuation. A finite bivariate degree truncation
is therefore \emph{not} an approximation uniformly valid below a single Hahn
valuation threshold. The accompanying computation checks a finite coefficient
identity, not a truncation of all coefficients below $1+\omega$.
\part*{Part III. Effectivity, verification, and boundaries}
\addcontentsline{toc}{section}{\textbf{Part III. Effectivity, verification, and boundaries}}

\section{How to compute the invariants, and what computation cannot decide}
\label{nab:sec:effectivity}

\subsection{A support-aware procedure for Criterion P}

For a presented star-cover datum the mathematical procedure is:
\begin{enumerate}
\item at each puncture compute the normalized factor by
\eqref{nab:eq:factorrec1}--\eqref{nab:eq:factorrec2}, or exploit a finite
matrix identity such as \eqref{nab:eq:HeisenbergP} or
\eqref{nab:eq:chainP};
\item extract the union of the supports of the normalized \emph{polar}
coefficients --- not the supports of the raw entries, and not just the leading
valuations;
\item supply a proof of well-ordering, or a descending-sequence obstruction,
for that union. In the admissible case use
\eqref{nab:eq:globalrec}--\eqref{nab:eq:qrec} to construct a splitting with a
monoid support certificate.
\end{enumerate}
These steps give an exact criterion, not a decision algorithm for an arbitrary
black-box infinite family of functions. In general an input description must
support a proof of the required infinite order property: checking many
exponents or many punctures does not settle it, and \cref{nab:cor:bounded}
shows that a finite-puncture computation provably cannot.

\subsection{Finite truncations}

For finitely generated positive rational exponent monoids every bounded
exponent interval is finite, so a sparse dictionary keyed by exact rational
exponents suffices for a finite-cutoff implementation of the local recursion,
each dictionary value being a matrix of ordinary Laurent polynomials, with
$\pi_-$ keeping the negative powers of $u$ and $\pi_+$ the nonnegative ones.
For a general higher-rank ordered group a bounded interval can be infinite and
no such finite-cutoff assumption is justified: an implementation must either
work with a certified finite dependency set or restrict its representation
class. This is the computational counterpart of \cref{nab:rem:nonarch}, and
\cref{nab:ex:boundedsupport} is a case where even the admissible answer cannot
be enumerated below a fixed bound in finitely many steps.

\subsection{What the Criterion M recursion actually depends on}

\begin{proposition}[Finite exponent dependence]\label{nab:prop:dependence}
Fix an admissible support $E$, a period section $\sigma$, a target exponent
$\gamma\in\angles E\setminus\{0\}$, and the meridian paths. The coefficient
$\Omega^\sigma_\gamma$ is determined by finitely many exponent-indexed vectors
of target coefficients $(m_{d,\delta})_{d\in D}$, $\delta\in E$, and finitely
many exponent patterns of ordinary products, iterated integrals, and
applications of $\sigma$. The set of punctures inside one such vector may still
be infinite.
\end{proposition}

\begin{proof}
Every lower coefficient entering the defect at $\gamma$ has a weight occurring
as a part in a decomposition of $\gamma$ into members of $\angles E$, and there
are finitely many such decompositions and parts by
\cref{nab:lem:finiteclosure}. Recursively replacing a constructed coefficient
by its own defining defect only refines one of these finite exponent-word
decompositions; equivalently all expressions are organized by finite rooted
trees whose leaves are $E$-letters summing to $\gamma$, and there are finitely
many possible leaf words by \cref{nab:lem:neumann} and finitely many tree
shapes and partitions for each fixed finite number of leaves. Applying $\sigma$
may use an entire element of $\C^D$, so this argument says nothing about
finite dependence on puncture coordinates.
\end{proof}

A relative algorithm is therefore explicit: represent the support well enough
to list the finitely many words relevant to a requested exponent, compute the
lower-weight iterated-integral prediction, subtract it from the target
coefficient vector, and apply the chosen ordinary period section. This does not
imply that the steps are computable for every set-theoretically specified Hahn
series, nor that a non-effective support can be enumerated, nor that an
infinite product vector can be supplied to a finite program. For finite
puncture sets the period section is explicit and the remaining tasks are
ordinary symbolic or numerical integrations.

\begin{proposition}[Compatibility with enlargement of the workspace]
\label{nab:prop:basechange}
Let $\iota:\Gamma\hookrightarrow\Gamma'$ be an order-preserving group
embedding, with the same complex coefficient data, meridians and period
section. For an admissible representation, its normalized realizer in the
larger Hahn workspace is obtained from the original realizer by replacing
$t^\gamma$ with $t^{\iota(\gamma)}$, and the same compatibility holds for
normalized fundamental matrices and framed comparison gauges.
\end{proposition}

\begin{proof}
Order embeddings preserve Hahn summability and addition of exponents, and the
finite coefficient formulas for products, iterated integrals and ordinary
primitives are unchanged, so the transfinite recursion on $\angles E$ maps to
the same recursion on $\angles{\iota(E)}$ with no extra coefficients outside
that monoid. Alternatively the transported family satisfies all defining
equations and normalizations, so uniqueness in
\cref{nab:thm:fundamental,nab:thm:RH,nab:thm:gauge} identifies it with the
larger-workspace construction.
\end{proof}

\subsection{Why an infinite linear period section cannot be continuous}

The choice-dependent section of \cref{nab:lem:period} is not merely missing an
estimate that could be supplied later. Give $\C^D$ its product topology and
give holomorphic one-forms on $U$ the topology of uniform convergence of the
coefficient functions on compact subsets, restricted to $\LL(D)$. The
following is an auxiliary observation, not a claim of novelty for the general
functional-analytic phenomenon.

\begin{proposition}[No continuous linear period right inverse]
\label{nab:prop:discontinuous}
If $D$ is infinite there is no continuous complex-linear right inverse
$\sigma:\C^D\to\LL(D)$ to $P$ for these topologies. No such continuous linear
right inverse exists even if its values are allowed to be arbitrary
holomorphic one-forms on $U$ with the specified meridian periods.
\end{proposition}

\begin{proof}
Choose a closed ordinary disk $K\subset U$ with nonempty interior, write a
one-form on $U$ as $f(z)\,dz$ and put $q_K(f\,dz)=\sup_{z\in K}|f(z)|$. If
$\sigma$ were continuous and linear, the continuous seminorm $q_K\circ\sigma$
on the product $\C^D$ would be bounded by a constant times a maximum of
finitely many coordinate absolute values, so there would be a finite
$F\subset D$ and $C_0>0$ with
\[
 q_K(\sigma(v))\leq C_0\max_{d\in F}|v_d|\qquad(v\in\C^D);
\]
this standard product-topology bound also follows by taking a basic
finite-coordinate neighbourhood on which the seminorm is less than one and
scaling. Choose $d\in D\setminus F$ and take the coordinate vector $e_d$. The
bound implies $q_K(\sigma(e_d))=0$, so the coefficient function of
$\sigma(e_d)$ vanishes on a disk and the ordinary identity theorem on the
connected domain $U$ makes it identically zero; its period around $\ell_d$ is
then zero, contradicting $(P\sigma(e_d))_d=1$.
\end{proof}

So the algebraic splitting used in the infinite-puncture normal form cannot
have product-topology continuity. This does not contradict
\cref{nab:thm:RH}, whose coefficient category imposes no continuity, and it
does not rule out nonlinear choices with other regularity properties, or useful
linear choices on smaller weighted sequence spaces; those are different
problems. For finite $D$ the explicit section \eqref{nab:eq:finitesection} is
continuous and the obstruction disappears.

\section{The accompanying exact and numerical checks}
\label{nab:sec:verification}

Two independent check suites were delivered with the two source manuscripts and
are preserved unchanged in \texttt{code/} and \texttt{data/} of this report
directory. Neither suite verifies a theorem; both audit finite identities, and
both print their own scope disclaimer. The facts recorded in this section were
confirmed by an independent run of the code, described in
\cref{nab:app:repro}.

\paragraph{Part~I: exact symbolic algebra.}
The script \texttt{08-polar-support-matrix-cousin-verify.py} uses Python with
SymPy and exact rational arithmetic, with no floating-point comparison in any
checked identity. It verifies: the rank-three factors and the missing cross
term, including $P H=G$, $\det=1$, nilpotence of $G-I_3$, the repaired product
of \cref{nab:ex:repair}, the matrix logarithm of \cref{nab:prop:logdiag} with
its $(1,3)$ coefficient $-xy/2u$ and the exact discrepancy $xy/2u$ between the
naive exponentiated-logarithm factor and the true polar factor; the chain
inverse, chain factorization, row-one square-zero property and determinant for
ranks $3$ through $10$ with symbolic $c_j$, together with the closed exponent
formulas \eqref{nab:eq:lowerexponents}--\eqref{nab:eq:lastexponent} and their
monotonicity in $n$ for ranks $4$ through $12$ and $n=2,\ldots,25$; a genuinely
noncommuting $2\times2$ series over the monoid generated by $\{1/2,2/3\}$
truncated at cutoff $3$ (sixteen exponents including zero), on which it runs
the actual $\pi_-/\pi_+$ recursion
\eqref{nab:eq:factorrec1}--\eqref{nab:eq:factorrec2} and checks $PH=G$
coefficientwise, that polar outputs are negative Laurent and regular outputs
holomorphic, holomorphic right-gauge invariance
$\Pol(GQ)=\Pol(G)$ with $\Reg(GQ)=\Reg(G)Q$, and a two-sided inverse; and an
exact four-puncture splitting with extension of every local frame. The
recorded run reports \textbf{2{,}899 exact checks passed} under Python 3.13.5
and SymPy 1.14.0; the independent rerun reproduced the same four PASS lines and
the identical count of 2{,}899 under Python 3.14.4 and SymPy 1.14.0, the
version banner being the only differing output line.

\paragraph{Part~II: exact algebra plus a coefficient-ODE audit.}
The script
\nolinkurl{09-support-monodromy-realization-verify-examples.py} has two
components. SymPy checks eight exact algebraic assertions:
$N_+^2=N_-^2=0$, $[N_+,N_-]=H$, both target determinants equal one, the exact
group commutator \eqref{nab:eq:commutator} and its determinant, and
tracelessness of the correction. A second component integrates the
degree-graded coefficient ODE system over the bivariate degrees
$(0,0),(1,0),(0,1),(2,0),(1,1),(0,2)$ with \texttt{scipy} \texttt{solve\_ivp}
\texttt{DOP853} at \texttt{rtol}~$2\cdot10^{-12}$ and
\texttt{atol}~$2\cdot10^{-13}$, along explicit meridian paths: from $1/2$ to
$1/4$, once positively around the radius-$1/4$ circle about $0$, and back; and
from $1/2$ to $3/4$, around the radius-$1/4$ circle about $1$ from angle $\pi$
to $3\pi$, and back. It does \emph{not} substitute floating-point surrogates
for $u$ or $v$. The uncorrected mixed term is compared with $\pm CH$ and the
corrected transport with the target through degree two. Recorded maximum
entrywise errors are $1.17\cdot10^{-13}$ and $1.47\cdot10^{-13}$ respectively
at both meridians; the assertions use a looser absolute tolerance of
$2\cdot10^{-9}$ so as not to depend on platform-specific last digits. The
recorded environment is Python 3.13.5 on Linux with NumPy 2.3.5, SciPy 1.17.0
and SymPy 1.14.0.

\paragraph{A warning about rerunning the Part~II script.}
That script \emph{writes its own evidence}: it emits
\texttt{verification\_results.json} beside itself, via
\texttt{Path(\_\_file\_\_).with\_name}. In the original archive layout this
could overwrite the recorded output beside the script. In the maintained
layout it creates or overwrites an output under \texttt{code/}, while the
separately named record under \texttt{data/} is unchanged. Reruns should use a
copy outside the report directory, as in \cref{nab:app:repro}; the delivered
\texttt{09-support-monodromy-realization-verification-results.json} in
\texttt{data/} is the original record.

\paragraph{What the checks do not establish.}
The Part~I checks audit algebraic signs, multiplication order and finite
coefficient identities; they do not prove the support lemma, Cartan's theorem
B, any infinite descending obstruction, an analytic existence theorem, or
novelty. The Part~II numerical results are observed floating-point errors, not
rigorous enclosures, and test only a finite bivariate degree-two identity; the
exact checks are algebraic identities in two formal variables. Neither suite
tests an arbitrary Hahn field, justifies a transfinite recursion, proves any
support set well ordered, verifies anything in a proof assistant, or
establishes that the results are unpublished. No implication in this report
depends on the outcome of either script: the general proofs are the arguments
in the body.

\section{Exactly what is not claimed}\label{nab:sec:scope}

The following list merges the honest limitations stated by both source
manuscripts. Every one of them survives this merge, and none is weakened.
Items \ref{nab:nc:refereed} through \ref{nab:nc:formal} come from Part~I's
scope, items \ref{nab:nc:positivity} through \ref{nab:nc:mueller} from
Part~II's, and the boundary note at the end records the one place where the two
parts make different standing assumptions.

\begin{enumerate}
\item\label{nab:nc:refereed}
Nothing here is refereed or proof-assistant verified. Both source manuscripts
are self-described AI-assisted research drafts and say so on their title pages
and in their READMEs.

\item\label{nab:nc:localization}
Nontriviality in Part~I is proved \emph{only} over the integral sheaf $\II$.
We explicitly do not claim that the examples survive after all monomials are
inverted, that is over $\HH$, because positive monomials become units there,
no reduction homomorphism $\HH\to\OO$ exists, and the framing-normalization
argument of \cref{nab:prop:integraltrivial} cannot be repeated. This is listed
as an open question, not as a believed consequence. In particular neither
criterion of this report is restated over a coefficient field in which
negative-valuation changes of frame are admitted.

\item\label{nab:nc:starcover}
The exact criterion of \cref{nab:thm:criterion} is only for the
puncture-and-disk (star) cover, where each overlap has a distinguished isolated
Laurent singularity. For general covers only the \emph{sufficient} globally
support-admissible theorem \cref{nab:thm:Stein} is available, and its
hypothesis \eqref{nab:eq:Ecover} is a condition on a \emph{presentation}, not
an invariant necessary condition: regular local gauges can destroy it without
destroying triviality.

\item\label{nab:nc:blackbox}
Neither criterion is a decision algorithm for a black-box infinite family. An
input description must support a proof of the infinite order property;
checking many exponents or many punctures does not settle it, and
\cref{nab:cor:bounded} shows that finite-puncture computation provably cannot
detect the Part~I examples.

\item\label{nab:nc:truncation}
Finite-cutoff implementation is justified only for finitely generated positive
rational exponent monoids. For a general higher-rank ordered group a bounded
exponent interval can be infinite.

\item\label{nab:nc:classification}
Part~I decides triviality rather than classifying all bundle isomorphism
classes. \Cref{nab:prop:torsor} describes all solutions once a splitting
exists, not the nonsplit objects. Part~II does classify framed positive gauge
classes by their admissible monodromy representations
(\cref{nab:thm:gauge}).

\item\label{nab:nc:higherdim}
The sharp star-cover invariant does not generalize to positive-dimensional
analytic singular loci: there is no one-variable Laurent projection there, and
an ideal-theoretic or cohomological replacement may have nontrivial relations.

\item\label{nab:nc:credits}
Credited as established mathematics and not new: the Neumann support calculus;
the ordinary Laurent, Mittag--Leffler, Cartan~B and Behnke--Stein inputs; and
algebraic noncommutative Birkhoff factorization in complete filtered algebras.
The local factorization is reproved only for explicit arbitrary-ordered-group
support control. The citation of \cite{EFGK} is to the \emph{algebra} of
recursive Birkhoff factorization and to the necessity of commutator
corrections; nothing in this report imports, endorses or reinterprets any
application of that algebra, and no assertion here is physical, empirical or
applied in any way. Likewise ``positive'' never means Hermitian positivity,
and the normalized polar factorization is \emph{not} the Hilbert-space polar
decomposition.

\item\label{nab:nc:priority}
The literature checks were targeted and are explicitly not certifications of
priority. Broad searches that returned irrelevant material were not treated as
evidence of absence, a repository code-search response marked incomplete was
not used as evidence of absence, and no named published conjecture is claimed
solved.

\item\label{nab:nc:code}
The companion code checks exact finite algebra and one finite numerical
example only; see the final paragraph of \cref{nab:sec:verification}.

\item\label{nab:nc:formal}
Formal verification is future work. The local support recursion, the chain
identity and the necessity directions appear well suited to formalization
with Hahn series and finite matrix algebra, but the global sufficiency proofs
additionally need a formal ordinary Mittag--Leffler theorem or a suitable
sheaf-cohomology interface. Neither script is a substitute for such a
formalization.

\item\label{nab:nc:positivity}
Positivity of exponents is essential throughout Part~II. There is no treatment
of nontrivial exponent-zero monodromy, of arbitrary ordinary residues with
resonances, or of arbitrary algebraic groups.

\item\label{nab:nc:notclassical}
\Cref{nab:thm:RH} is not an unrestricted classical Riemann--Hilbert theorem and
not a Riemann--Hilbert correspondence on a general surcomplex domain; likewise
\cref{nab:thm:criterion} is not the unrestricted classical Cousin theorem over
$\HH$. The base, the universal cover and the path category are \emph{ordinary}
complex by design, not provisional notation for a fine-topological
construction.

\item\label{nab:nc:finetop}
There is no fine-topological theory of surcomplex loops here. Convergence of
Hahn partial sums in the fine topology is never asserted, and
\cref{nab:ex:boundedsupport} shows that partial sums can fail to converge even
in the valuation topology. The surcomplex reading is an interpretation of a
Hahn model via $t^\gamma=\omega^{-\gamma}$ and does not assert an ambient
surcomplex manifold.

\item\label{nab:nc:choice}
The infinite-puncture period section is obtained by a vector-space splitting
using the axiom of choice. It is not continuous ---
\cref{nab:prop:discontinuous} proves that impossible. No norm estimate or
effective procedure is supplied; only an ordinary complex-linear splitting is
claimed, with no computability assertion about a specified representation of
the infinite input data.

\item\label{nab:nc:infinity}
For infinitely many punctures, no regularity or growth condition at infinity is
imposed on realizers. This is deliberate, and the freedom to use arbitrary
ordinary entire correction terms at each exponent is important to the existence
proof.

\item\label{nab:nc:nosheaf}
Part~II does not assume that the global-support assignment is a sheaf for
arbitrary covers, uses no Cartan theorem and no sheaf cohomology for a Hahn
coefficient sheaf, and relies on no companion claim about Hahn Noetherianity or
a Hahn Nullstellensatz.

\item\label{nab:nc:deformations}
Deformations of a \emph{nontrivial ordinary connection} are open. The
linearized operator there becomes a twisted period map with values in an
adjoint local system rather than the scalar period map $P$, and the triangular
proof of \cref{nab:lem:triangular} supplies no right inverse for it. This is
not to be confused with \cref{nab:thm:deformation}, which deforms a nontrivial
ordinary \emph{vector bundle} on a Stein base and does have its cohomological
input available.

\item\label{nab:nc:dimension}
Higher-dimensional bases are open for Part~II: flatness is automatic only for
holomorphic one-forms on a curve, a proposed correction must satisfy the
nonlinear flatness equation as well as the period conditions, and the result
cannot be inferred by changing $z$ to a tuple.

\item\label{nab:nc:numerics}
The Part~II numerical checks are observed floating-point errors, not rigorous
enclosures, and test only a finite bivariate degree-two identity.

\item\label{nab:nc:mueller}
The Hahn-meromorphic category of M\"uller--Strohmaier \cite{Muller} is
explicitly \emph{not} identified with ours: there the parameter is tied to an
expansion near zero on a logarithmic cover, whereas here $t$ is an independent
Hahn parameter and $z$ the ordinary analytic variable.

\item\label{nab:nc:common}
No theorem of either part is claimed to be an instance of a common
generalization proved here. The two criteria are two theorems with two proofs
(\cref{nab:sec:twocriteria}); they share only
\cref{nab:lem:neumann,nab:lem:inverse} and an ordinary Mittag--Leffler input.
\end{enumerate}

\paragraph{Boundary note on the sheaf convention.}
As recorded in \cref{nab:sub:conventions}, Part~I proves and uses the sheaf
property of the integral common-domain assignment on the ordinary base, while
Part~II deliberately assumes nothing of the kind. Both are correct as stated,
the difference is one of standing hypotheses rather than of fact, and no
theorem of Part~II is strengthened here by importing Part~I's convention.

\section{Questions exposed by the results}\label{nab:sec:open}

\paragraph{Monomial localization.}
Does $\mathcal E_G\otimes_{\II}\HH$ remain nontrivial for the explicit hidden
and deep examples? A proof would need control of general Hahn matrix gauges,
not just positive ones, and the ordinary reduction normalization is unavailable
after localization. This question is intentionally not answered.

\paragraph{Arbitrary-cover normal forms.}
The polar support is exact on the puncture-and-disk cover because every overlap
has a distinguished isolated Laurent singularity. An equally explicit invariant
for general covers would require a substitute for that canonical local negative
factor; \cref{nab:thm:Stein} does not provide one.

\paragraph{Classification beyond triviality.}
Deciding whether a positive star-cover cocycle represents the trivial integral
bundle is not classifying all isomorphism classes, which would require
understanding the action of global and local gauges on \emph{pairs} of polar
families.

\paragraph{A common generalization of the two criteria.}
Both criteria say that a union of supports must be well ordered, over different
objects (\cref{nab:sec:twocriteria}). Is there a single statement, in some
category of positive Hahn objects on a punctured base with a distinguished
local normalization, of which both are instances? We do not know, and nothing
in this report should be read as evidence that the naive common statement ---
taking the union over raw local data --- is true: \cref{nab:thm:hidden} and
\cref{nab:ex:repair} refute it for Criterion P.

\paragraph{Deformations of a nontrivial ordinary connection.}
Fix an ordinary connection $\Omega_0$ and a representation reducing to its
monodromy rather than to the identity. The first correction operator is then a
twisted period map into the adjoint local system. A precise extension should
specify when that map has a right inverse within a chosen pole class, and
should track local resonances.

\paragraph{Controlled growth and effective infinite data.}
For a prescribed weight or growth condition on residues at $D$, identify
spaces on which a useful continuous or computable period section exists. The
exact question is to characterize the residue sequence spaces admitted by the
chosen coefficient growth class, and then to prove that nonlinear monodromy
corrections stay inside them; arbitrary algebraic splittings need not preserve
any such structure.

\paragraph{Higher-dimensional bases.}
For Stein manifolds of dimension greater than one the support-controlled
rigidity theorems still apply, but the sharp star-cover invariant does not
generalize, and for connections the nonlinear flatness equation enters. Both
directions are open.

\paragraph{Formal verification.}
See item~\ref{nab:nc:formal} of \cref{nab:sec:scope}.

\section{Conclusion}\label{nab:sec:conclusion}

In the positive integral Hahn category over a set-sized ordered value group,
both the multiplicative gluing problem and the inverse monodromy problem have
exact answers, and in both the answer is a well-ordering condition on a union
of Hahn supports. They are not the same condition.

For gluing, the correct obstruction is not the support of the input entries,
nor the support of their raw principal parts, nor a determinant. One must first
remove the holomorphic factors in the correct noncommutative order; the
resulting normalized polar support is then an exact necessary and sufficient
criterion on a puncture-and-disk cover. The explicit examples show why that
normalization is consequential: a family with one singular input exponent can
force a descending support at a hidden matrix entry; adding a cross term can
repair the family while worsening its raw singular support; and the same
phenomenon can be postponed to the final central layer of upper-unitriangular
groups of arbitrarily large finite rank, yielding nontrivial determinant-one
integral Hahn bundles with trivial ordinary reduction that are trivial on every
bounded plane domain. At the opposite boundary, one globally well-ordered
support certificate restores Stein rigidity in every finite rank; cyclic value
groups supply such a certificate automatically and noncyclic ones do not.

For monodromy, the correct obstruction is failure of well-ordering of the raw
support union of the prescribed matrices. This condition is an invariant of
the representation. Realizers
can be taken coefficientwise logarithmic, a fixed linear period section
selects a unique normalized one, framed positive gauges classify connections
with equal monodromy, and within the logarithmic class the comparison gauge
extends across every puncture. The order of the value group may have several
ranks, the singularity set may be infinite, and the monodromy may be
noncommutative; none of this requires a false topological limiting argument.

What both halves separate cleanly is ordinary analytic solvability,
noncommutative compatibility, and generalized-series admissibility. The
analytic work is done coefficientwise on the ordinary base; the Hahn support
calculus is exactly what makes those ordinary constructions coexist in one
surcomplex workspace; and the noncommutative correction --- the cross term
$-a_-c_+$ in Part~I, the mixed commutator $-uv\,CH(\eta_0-\eta_1)$ in Part~II
--- is what the abelianization cannot see.
\appendix
\crefalias{section}{appendix}
\crefalias{subsection}{appendix}

\section{A proof and support audit}\label{nab:app:audit}

The table records every place where an infinite formal expression occurs,
together with the set that controls it and the ordinary analytic fact required.
A support in the middle column is a \emph{containing} support; cancellation may
make the actual support smaller.

\medskip\noindent\textbf{Part~I constructions.}
\begin{center}\small
\begin{tabular}{@{}L{0.26\textwidth}L{0.35\textwidth}L{0.29\textwidth}@{}}
\toprule
Construction & Support control & Ordinary analytic requirement\\
\midrule
Positive matrix inverse, $\exp$, $\log$ &
Positive monoid generated by the input support; finitely many exponent words
and finitely many matrix paths at each coefficient
(\cref{nab:lem:neumann,nab:lem:inverse}) &
Coefficients remain on the original domain.\\ \addlinespace[0.35em]
Local polar factorization &
$\angles{\supp(G-I_r)}$; well-founded triangular recursion &
Laurent splitting on one fixed punctured disk, with no shrinking
(\cref{nab:lem:laurent}).\\ \addlinespace[0.35em]
Necessity of Criterion P &
All $P_a-I_r$ lie in the one monoid generated by $\supp(B_0-I_r)$ &
Right-gauge invariance and uniqueness of the local factorization.\\ \addlinespace[0.35em]
Sufficiency of Criterion P &
$\angles{\EU(G)}$, for the central matrix and for the $Q_a$ &
One ordinary isolated-singularity interpolation per matrix coefficient and
exponent (\cref{nab:lem:mlplane,nab:lem:mlsurface}).\\ \addlinespace[0.35em]
Disk trivializers &
Generated monoid of the global polar support together with that disk's own
original support &
Products of matrices holomorphic on that disk.\\ \addlinespace[0.35em]
Right torsor of splittings &
Finite union of the two solution supports, then positive monoid closure &
Identity-theorem extension across each puncture.\\ \addlinespace[0.35em]
Stein splitting &
One monoid generated by all transition supports &
Ordinary $H^1(X,\OO^{r^2})=0$.\\ \addlinespace[0.35em]
Deformation rigidity &
The same positive monoid; ordinary transitions are exponent-zero coefficients
&
Ordinary $H^1(X,\Endsh(\mathcal V))=0$.\\ \addlinespace[0.35em]
Rank-three obstruction &
Forced descending set $\{1+1/n\}$ &
Only simple poles are used.\\ \addlinespace[0.35em]
Depth-$(r-1)$ obstruction &
Lower partial-sum supports well ordered; final support $\{r-2+1/n\}$ descends
&
Only simple poles and holomorphic constants are used.\\ \addlinespace[0.35em]
Integral nontriviality &
A hypothetical integral frame normalizes to a positive frame &
Reduction of invertible matrices and gluing of the reduced frame.\\ \addlinespace[0.35em]
\bottomrule
\end{tabular}
\end{center}

\medskip\noindent\textbf{Part~II constructions.}
\begin{center}\small
\begin{tabular}{@{}L{0.26\textwidth}L{0.35\textwidth}L{0.29\textwidth}@{}}
\toprule
Construction & Support control & Ordinary analytic requirement\\
\midrule
Fundamental matrix on the cover &
$\angles{\supp\Omega}$, for $Y$ and for $Y^{-1}$; the same bound for
\emph{every} holonomy matrix &
Holomorphic primitives on a simply connected Riemann surface.\\ \addlinespace[0.35em]
Dyson transport &
Strong summability by the finite-word clause; each coefficient a finite sum of
iterated integrals &
Ordinary iterated path integrals along an ordinary path.\\ \addlinespace[0.35em]
Necessity of Criterion M &
$E_\rho\subseteq\angles{\supp\Omega}\setminus\{0\}$, uniformly in the group
element &
None beyond the fundamental-matrix theorem.\\ \addlinespace[0.35em]
Sufficiency of Criterion M &
$\angles{E_\rho}$; transfinite recursion with no limit-ordinal passage &
Ordinary Mittag--Leffler through the period map $P$, plus one fixed linear
section $\sigma$ (axiom of choice for infinite $D$).\\ \addlinespace[0.35em]
Uniqueness of the normalized realizer &
Least differing exponent in a well-ordered union &
Injectivity of $P$ on $\im\sigma$.\\ \addlinespace[0.35em]
Framed gauge and its inverse &
$\angles{\supp\Omega\cup\supp\Psi}$ &
Single-valuedness on $U$; descent of coefficients from the cover.\\ \addlinespace[0.35em]
Entire extension of a logarithmic gauge &
The same monoid, unchanged by the extension &
A simple pole with zero residue is removable.\\ \addlinespace[0.35em]
Positive Frobenius form &
$\angles E$ for $H$ and $H^{-1}$ &
Primitives on a disk; $H_\beta(0)=0$ makes $H_\beta/z$ holomorphic.\\ \addlinespace[0.35em]
Traceless refinement &
No change of support &
Finite polynomial determinant identity; one entrywise linear section.\\ \addlinespace[0.35em]
Finite exponent dependence &
Finitely many decompositions and parts inside $\angles E$
(\cref{nab:lem:finiteclosure}) &
Finitely many rooted trees and ordinary operations per exponent.\\
\bottomrule
\end{tabular}
\end{center}

The indispensable hypotheses are these. The value group is a set and is
totally ordered. Supports are well ordered, not merely bounded below. Matrix
rank is finite. The puncture set is ordinary closed and discrete. Coefficient
functions share their stated ordinary domains. In Part~I the gluing matrices
reduce to the identity, or to a specified ordinary bundle in
\cref{nab:thm:deformation}; in Part~II the representation is near-identity with
positive support. None of these hypotheses is replaced anywhere by a
fine-topological convergence assumption.

\section{A dependency map}\label{nab:app:dep}

The Part~I chain is
\[
\begin{gathered}
\text{Neumann support lemma}+\text{ordinary Laurent splitting}\\
\Downarrow\\
\text{local normalized factorization and right-gauge invariance}\\
\Downarrow\quad\text{with ordinary isolated-singularity interpolation}\\
\text{exact global polar-support criterion (Criterion P)}\\
\Downarrow\\
\text{rank-three cross-term obstruction and arbitrary-depth examples}\\
\Downarrow\quad\text{with reduction of an integral frame}\\
\text{nontrivial integral vector bundles.}
\end{gathered}
\]
The arbitrary-cover Stein theorem is a separate positive branch whose only
cohomological input is ordinary coherent-sheaf acyclicity. It is not used to
justify the infinite star-cover criterion on the plane, whose interpolation
step has its own direct proof; it is used for the statement about all bounded
restrictions and for the universal cyclic-group result.

The Part~II chain is
\[
\begin{gathered}
\text{Neumann support lemma}\\
\Downarrow\\
\text{support control for inverses and Dyson expansions}\\
\Downarrow\quad\text{with ordinary primitives on a simply connected cover}\\
\text{fundamental matrix and the triangular variation lemma}\\
\Downarrow\quad\text{meeting ordinary Mittag--Leffler through }P\text{ and }
\sigma\\
\text{exact realization recursion (Criterion M)}\\
\Downarrow\\
\text{framed gauge classification, entire extension, }\SL_r\text{ refinement.}
\end{gathered}
\]
The gauge classification uses only the fundamental matrices; its entire
extension uses only single-valuedness and removability of a simple pole with
zero residue. The traceless refinement adds a finite polynomial determinant
identity. The counterexamples use the necessity half of
\cref{nab:thm:RH} and explicit support orderings. No implication in either
chain depends on the outcome of a verification script.

The two chains meet only at \cref{nab:lem:neumann,nab:lem:inverse} and at the
shared observation that the ordinary analytic input is applied one exponent at
a time. Neither criterion is derived from the other.

\section{Provenance, reproduction, and the merge record}\label{nab:app:repro}

\subsection{Sources of this report}

This report merges two independently written manuscripts, both dated
September 21, 2026, both delivered with their own code and recorded output, the
latter preserved unchanged in \texttt{code/} and \texttt{data/} of this
directory:

\begin{itemize}
\item \emph{Nonabelian Support Obstructions in Surcomplex Analysis: A Sharp
Matrix Cousin Criterion and Arbitrarily Deep Unipotent Examples}, from the
archive \texttt{surcomplex\_nonabelian\_support}; source file
\texttt{08-polar-support-matrix-cousin.tex}, delivered with its README. It
contributes Part~I entire, the shared coefficient category and support
calculus, the Laurent and isolated-singular-part inputs, the normalized polar
factorization, the exact Criterion P with its invariance corollaries and torsor
statement, the rank-one recovery, the rank-three obstruction with its repair
and logarithm diagnostic, the chain identity and depth hierarchy, the integral
vector-bundle statements with the localization boundary, and the Stein
splitting, deformation rigidity and cyclic dichotomy.

\item \emph{Global Hahn Support and the Surcomplex Riemann--Hilbert Problem:
Noncommutative monodromy, countably many singularities, and an exact
realization obstruction}, from the archive
\texttt{surcomplex\_riemann\_hilbert}; source file
\texttt{09-support-monodromy-realization.tex}, delivered with its README and
its
\texttt{BUILD\_AND\_PROVENANCE.json}. It contributes Part~II entire: strong
summability and the finite-decomposition lemma, the free-meridian lemma, the
period map with its linear section, the fundamental-matrix and triangular
lemmas, Criterion M with its intrinsic and conjugation-invariance forms, the
framed gauge classification and entire extension, the positive Frobenius form,
the finite Fuchsian specialization and the $\SL_r$ refinement, the three
calibrating examples, the mixed commutator correction with its $\Q+\Q\omega$
specialization, the finite-dependence and base-change propositions, and the
impossibility of a continuous linear period right inverse.
\end{itemize}

Where the two manuscripts proved the same fact --- the Neumann support calculus
as an input, and the positive inversion, exponential and logarithm lemma ---
the statement is printed once, in
\cref{nab:lem:neumann,nab:lem:inverse}, and both members verified it
independently in the same form. Where they supplied genuinely different routes
to a needed ingredient, both routes are kept and marked: the ordinary
Mittag--Leffler input appears twice, as the unbounded-order
isolated-singular-part construction
\cref{nab:lem:mlplane,nab:lem:mlsurface} and as the linear period section
\cref{nab:lem:period}, for the reasons given at the end of
\cref{nab:sec:ordinary}; and the impossibility of repair by enlarging the value
group is proved twice, for two different objects, in
\cref{nab:cor:extension} and \cref{nab:prop:conjugation}. The local
factorizations of \cref{nab:thm:local} and \cref{nab:thm:frobenius} are
likewise two different normalizations of two different objects, and neither is
derived from the other here.

\subsection{Repository provenance}

Repository: \url{https://github.com/VladimirReshetnikov/Surreal}, reviewed at
commit \texttt{39f2be6667ade51bca2b45daa47e289d69c09764}, inspected
September 21, 2026. No repository files were changed by either source
manuscript and nothing was pushed by them. The report catalogue
\cite{RepoMap} supplies the report map and the distinctions between coefficient
rings, which both members call essential. The global-divisor report
\cite{RepoGlobal} is the source of the motivating open direction and of the
scalar results we recover rather than reprove; the differential-equations
report \cite{RepoDiff} supplies the forward monodromy map and the
ordinary-paths wording. The analysis, analytic-geometry and contour packages
\cite{RepoAnalysis,RepoGeometry,RepoContours} are credited by Part~II as
antecedents for scalar global support obstructions and for the coefficientwise
differential and integration viewpoint; they are not contributions claimed
here, and no draft ring theorem or Hahn sheaf-cohomology theorem from them is
used in any proof.

Neither source manuscript conducted an exhaustive audit of every source
manuscript in the repository, and one repository code-search response was
marked incomplete; its empty result is not evidence of absence. The comparison
above is with the documented contents actually reviewed.

\subsection{Building this report}

The source is standalone, with an internal bibliography, no external
bibliography file, no graphics and no separately distributed fonts. A standard
TeX Live installation builds it with
\begin{verbatim}
latexmk -pdf -interaction=nonstopmode article.tex
\end{verbatim}
Alternatively, run \texttt{pdflatex} with
\texttt{-interaction=nonstopmode -halt-on-error} three times to settle the
cross-references and contents.

\subsection{Rerunning the delivered checks}

The two scripts are preserved under \texttt{code/} with their original
contents and their recorded outputs under \texttt{data/}. The Part~I script
needs Python 3.10 or later with SymPy; the Part~II script needs NumPy, SciPy
and SymPy. From this report directory, run copies in a temporary directory:
\begin{verbatim}
check_dir=$(mktemp -d)
cp code/*-verify*.py "$check_dir/"
cd "$check_dir"
python 08-polar-support-matrix-cousin-verify.py
python 09-support-monodromy-realization-verify-examples.py
\end{verbatim}

Two operational cautions, both verified directly rather than taken from a
README.

First, the Part~II script writes \texttt{verification\_results.json} beside
itself. An in-place run in the maintained layout writes under \texttt{code/};
in the original archive layout it could overwrite the historical record.
Running a temporary copy keeps generated evidence separate from both the
script and the delivered record.

Second, the Part~II README as delivered lists a file
\texttt{SHA256SUMS.txt} among the archive contents. That file is
\emph{not} present in the delivered directory and is not present here, so no
checksum verification of the original archive is possible from what was
shipped.

An independent rerun of both suites, on copies and not in the extraction tree,
found: the Part~I suite exits zero and prints its four PASS lines and
\texttt{SUCCESS: 2899 exact checks passed}, identical in count to the recorded
run, under Python 3.14.4 with SymPy 1.14.0 against the recorded Python 3.13.5
with SymPy 1.14.0, the version banner being the only differing line; and the
Part~II suite exits zero with all eight exact checks and all four numerical
error checks passing. The recorded Part~II build metadata claims zero warnings,
zero overfull or underfull boxes, zero undefined references, a 25-page
document, and a Poppler render review of all pages, and records
\texttt{formal\_verification: false}, \texttt{priority\_certified: false},
\texttt{peer\_reviewed: false}. Those three flags remain false for this merged
report.

\begin{thebibliography}{99}
\small\raggedright
\setlength{\itemsep}{0.3em}
\addcontentsline{toc}{section}{References}

\bibitem{RepoMap}
V.~Reshetnikov, \emph{Surreal repository: report catalogue},
\texttt{docs/README.md}, commit
\texttt{39f2be6667ade51bca2b45daa47e289d69c09764}, inspected September 21,
2026.
\url{https://github.com/VladimirReshetnikov/Surreal/blob/39f2be6667ade51bca2b45daa47e289d69c09764/docs/README.md}.

\bibitem{RepoGlobal}
\emph{Global Support Obstructions: Moving Divisors, Mittag--Leffler, and a
Picard Dichotomy}, AI-assisted merged research draft in V.~Reshetnikov's
\emph{Surreal} repository, \texttt{docs/surcomplex/global-divisors/}. Pinned
commit as above; see especially the exact Cousin criterion (Theorem~13.1), the
Picard-group vanishing dichotomy, and the subsection ``Further precise
directions''.
\url{https://github.com/VladimirReshetnikov/Surreal/blob/39f2be6667ade51bca2b45daa47e289d69c09764/docs/surcomplex/global-divisors/article.tex}.

\bibitem{RepoDiff}
\emph{Differential Algebra and Differential Equations over the Surreal and
Surcomplex Numbers}, AI-assisted merged research draft in V.~Reshetnikov's
\emph{Surreal} repository, \texttt{docs/surcomplex/differential-equations/}.
Pinned commit as above; see especially the theorem ``Monodromy valued in
$\GL_d(K_\Gamma)$'' and the ordinary-path remark following it.
\url{https://github.com/VladimirReshetnikov/Surreal/blob/39f2be6667ade51bca2b45daa47e289d69c09764/docs/surcomplex/differential-equations/article.tex}.

\bibitem{RepoAnalysis}
\emph{Surcomplex analysis: holomorphic functions on $\No[i]$},
\texttt{docs/surcomplex/analysis/README.md}, same repository and commit.
AI-assisted, unrefereed research documentation.

\bibitem{RepoGeometry}
\emph{Infinitesimal analytic geometry over the surcomplex numbers:
radius-free germs, a Nullstellensatz, and finite maps},
\texttt{docs/surcomplex/analytic-geometry/README.md}, same repository and
commit. AI-assisted, unrefereed research documentation.

\bibitem{RepoContours}
\emph{Contours and Stokes: Jordan separation, Cauchy theory, and Stokes
calculus on the surcomplex plane},
\texttt{docs/surcomplex/contours-and-stokes/README.md}, same repository and
commit. AI-assisted, unrefereed research documentation.

\bibitem{Neumann}
B.~H.~Neumann, On ordered division rings,
\emph{Transactions of the American Mathematical Society} \textbf{66} (1949),
202--252.
\url{https://www.jstor.org/stable/1990552}.

\bibitem{Poonen}
B.~Poonen, \emph{Maximally complete fields},
L'Enseignement Math\'ematique (2) \textbf{39} (1993), 87--106.
DOI: \url{https://doi.org/10.5169/seals-60414}.

\bibitem{vdDE}
L.~van den Dries and P.~Ehrlich, Fields of surreal numbers and exponentiation,
\emph{Fundamenta Mathematicae} \textbf{167} (2001), 173--188.
\url{https://doi.org/10.4064/fm167-2-3}.

\bibitem{MantovaMatusinski}
V.~Mantova and M.~Matusinski, \emph{Surreal numbers with derivation, Hardy
fields and transseries: a survey}, arXiv:1608.03413.
\url{https://arxiv.org/abs/1608.03413}.

\bibitem{Bournez}
O.~Bournez and Q.~Guilmant, \emph{Surreal fields stable under exponential and
logarithmic functions}, preprint (2022), arXiv:2201.08199. Cited for
normal-form and set-sized-surreal-field background only.
\url{https://arxiv.org/abs/2201.08199}.

\bibitem{EFGK}
K.~Ebrahimi-Fard, L.~Guo, and D.~Kreimer,
Spitzer's identity and the algebraic Birkhoff decomposition in pQFT,
\emph{Journal of Physics A: Mathematical and General} \textbf{37} (2004),
11037--11052. Cited solely as the closest algebraic antecedent for recursive
Birkhoff factorization in complete filtered noncommutative algebras and for the
necessity of commutator corrections; no application of that algebra is imported
here.
\url{https://arxiv.org/abs/hep-th/0407082}.

\bibitem{EFMP}
K.~Ebrahimi-Fard, D.~Manchon, and F.~Patras, \emph{A noncommutative
Bohnenblust--Spitzer identity for Rota--Baxter algebras solves Bogoliubov's
recursion}, arXiv:0705.1265.
\url{https://arxiv.org/abs/0705.1265}.

\bibitem{BS}
H.~Behnke and K.~Stein, Entwicklung analytischer Funktionen auf Riemannschen
Fl\"achen, \emph{Mathematische Annalen} \textbf{120} (1947/1949), 430--461.
\url{https://doi.org/10.1007/BF01447838}.

\bibitem{Cartan}
H.~Cartan, Vari\'et\'es analytiques complexes et cohomologie, in
\emph{Colloque sur les fonctions de plusieurs variables}, Brussels, 1953,
41--55. Classical source for coherent-sheaf acyclicity on Stein manifolds.

\bibitem{Schmitt}
M.~M.~Schmitt, \emph{Kohomologie mit Schranken und Fortsetzung holomorpher
Funktionen durch lineare stetige Operatoren}, doctoral thesis,
arXiv:math/0202263 (2002). Includes a proof of Cartan's theorem B for coherent
subsheaves of finite free holomorphic sheaves and coefficientwise coboundary
constructions.
\url{https://arxiv.org/abs/math/0202263}.

\bibitem{Wone}
O.~Wone, \emph{The Birkhoff--Grothendieck theorem}, arXiv:2304.08037,
version~2. Historical and mathematical exposition of bundle splitting and
Birkhoff factorization; its usual single-loop setting differs from the
infinitely many independently supported gluing data treated here.
\url{https://arxiv.org/abs/2304.08037}.

\bibitem{Chen}
K.-T.~Chen, \emph{Iterated path integrals},
Bulletin of the American Mathematical Society \textbf{83} (1977), 831--879.
DOI: \url{https://doi.org/10.1090/S0002-9904-1977-14320-6}.

\bibitem{Deligne}
P.~Deligne, \emph{\'Equations diff\'erentielles \`a points singuliers
r\'eguliers}, Lecture Notes in Mathematics, vol.~163, Springer, 1970. A
classical reference point, not a theorem reproved here in its original
generality.
DOI: \url{https://doi.org/10.1007/BFb0061194}.

\bibitem{Ilyashenko}
Y.~Ilyashenko and S.~Yakovenko, \emph{Lectures on Analytic Differential
Equations}, Graduate Studies in Mathematics, vol.~86, American Mathematical
Society, 2008.
DOI: \url{https://doi.org/10.1090/gsm/086}.

\bibitem{Muller}
J.~M\"uller and A.~Strohmaier, \emph{The theory of Hahn-meromorphic functions,
a holomorphic Fredholm theorem, and its applications},
Analysis \& PDE \textbf{7} (2014), no.~3, 745--770. Their analytic category is
explicitly not identified with ours; see item~\ref{nab:nc:mueller} of
\cref{nab:sec:scope}.
DOI: \url{https://doi.org/10.2140/apde.2014.7.745}.

\end{thebibliography}
\end{document}
